% !TeX program = lualatex
% !TeX encoding = utf8
% !TeX spellcheck = uk_UA
% !TeX root = ../main.tex

%=========================================================
\Opensolutionfile{answer}[\currfilebase/\currfilebase-Answers]
\Writetofile{answer}{\protect\section*{\nameref*{\currfilebase}}}
\chapter{Коливальні процеси в електричних колах}\label{\currfilebase}\hypertarget{EMOscillations}{}
\graphicspath{{\currfilebase/Pictures}}
%=========================================================
\epigraph{\Alexander Однакові рівняння мають однакові розв'язки.}{Р.~Фейнман, \enquote{Фейнманівські
лекції з фізики}, т.~2, гл.~12 (1964)}


Досі електричне і магнітне поле розглядалися переважно як самостійний фізичний об'єкт: ми записали рівняння Максвелла, з'ясували, що вони допускають
розв'язки у вигляді електромагнітних хвиль, і простежили, як збурення поля поширюється в просторі зі скінченною швидкістю $c$. Проте є практично важливі задачі, де поле не поширюється у вільному просторі, а виявляється \enquote{замкненим} в електричному колі --- системі провідників,
конденсаторів і котушок, з'єднаних певним чином.

Виникає природне питання: чи можна описувати процеси в такому колі, не розв'язуючи щоразу повні рівняння Максвелла, а обмежившись значно простішими
співвідношеннями --- законом Ома, правилами Кірхгофа, формулами для заряду конденсатора чи струму в котушці? Виявляється, що так, --- але лише за певної
умови, яка й дозволяє знехтувати хвильовим (запізненим) характером поля всередині кола. Ця умова називається \emphx{квазістаціонарним наближенням}\index{Квазістаціонарне наближення}, і саме
з неї ми розпочнемо вивчення коливальних процесів в електричних колах.

%% --------------------------------------------------------
\section{Квазістаціонарне наближення}\label{sec:quasi-stationary-approximation}\hypertarget{quasi-stationary-approximation}{}
%% --------------------------------------------------------

Процес вважається \emphx{квазістаціонарним}, якщо миттєві значення струму в усіх ділянках провідника (нерозгалуженого кола) однакові, а електричні поля в конденсаторах такі самі, як в електростатиці.

Сформулюємо умову квазістаціонарності. По електричному колу поширюється електромагнітний сигнал, швидкість якого порядку швидкості світла. Якщо частота сигналу $\omega$, то відповідна довжина хвилі порядку $\lambda = 2\pi c/\omega = cT$, де $T = 2\pi/\omega$ --- характерний період коливань або змін сигналу. Якщо довжина кола $l$ така, що
\begin{equation}\label{eq:quasistationary_condition}
	l \ll \lambda, \qquad \text{або} \qquad l \ll cT,
\end{equation}
то в усіх ділянках кола струми й напруги змінюються синхронно, в одній і тій самій фазі. Оскільки час проходження сигналом усього кола $\tau \sim l/c$, то умову квазістаціонарності можна записати у вигляді
\begin{equation}\label{eq:quasistationary_time}
	\tau \ll T.
\end{equation}

\begin{Attention}
	Наскільки серйозна ця умова для звичної нам мережі 50 Гц? Довжина хвилі становить $\lambda = cT \approx 3\times10^{8}\ \text{м/с}\cdot 0{,}02\ \text{с} \approx 6000$ км. Для будинкової проводки чи навіть цілого заводу це запас у сотні тисяч разів --- квазістаціонарність виконується з величезним надлишком. Але магістральна лінія електропередачі (ЛЕП) може тягнутися на сотні й тисячі кілометрів, тобто вже становити відчутну частку довжини хвилі мережі. Тоді просту схему з $R$, $L$, $C$ доводиться замінювати повноцінною теорією довгих ліній, в якій напруга і струм у різних точках проводу зсунуті за фазою, як у хвилі, що біжить уздовж лінії. Саме тому надто довга лінія, розімкнута на одному кінці, може несподівано увійти в резонанс із частотою мережі --- про це ще буде мова далі у цьому розділі.
\end{Attention}

Далі ми розглядаємо саме квазістаціонарні процеси. Це означає, що ми нехтуємо скінченним часом поширення електромагнітної взаємодії вздовж кола й вважаємо, що струм і напруга встигають узгодитися по всьому колу в кожний момент часу. Іншими словами, коло описується не як система з розподіленими параметрами (довга лінія), а як система із зосередженими елементами --- резисторами $R$, котушками $L$ і конденсаторами $C$, з'єднаними ідеальними провідниками. Саме в межах цього наближення справедливі закон Ома, правила Кірхгофа та рівняння коливального контуру, тож усі подальші результати --- вільні й вимушені коливання, змінний струм, резонанс --- чинні лише тією мірою, якою виконується умова \eqref{eq:quasistationary_time}.

%% --------------------------------------------------------
\section{Вільні коливання в коливальному контурі}\hypertarget{free-oscillations}{}
%% --------------------------------------------------------

%% --------------------------------------------------------
\subsection*{Рівняння коливального контуру}
%% --------------------------------------------------------

Нехай\index{Коливальний контур} є електричне коло, що включає ємність (конденсатор) $C$, індуктивність $L$, опір $R$ і стороннє джерело електрорушійної сили (ЕРС) $\emf$ (рис.~\ref{tikz:LCR_circuit}). Складемо рівняння, що визначає зміну струму в колі\footnote{У цьому розділі всюди ми використовуємо систему одиниць СІ, в якій $\Phi = LI$, $\emf = -d\Phi/dt$.}. Застосуємо теорему про циркуляцію електричного поля до замкнутого контуру $\Gamma = \{12341\}$:
\begin{equation}\label{eq:circulation}
	\oint\limits_{\{12341\}} E_l \, dl = -\frac{d\Phi}{dt}.
\end{equation}

В праву частину цієї рівності входить ЕРС індукції, що виникає в
котушці індуктивності при зміні магнітного потоку. Додатковий
облік цього елемента кола в лівій частині рівності вже не
потрібен. Будемо вважати також, що весь опір кола об'єднано в
один елемент $R$. Виберемо за додатний напрямок обходу контуру
той, що вказано на рис.~\ref{tikz:LCR_circuit} для струму $I$.

\begin{figure}[h!]
	\centering
	\localinput{LCR_circuit.tikz}
	\caption{Контур, що включає конденсатор, індуктивність, опір $R$ і стороннє джерело ЕРС $\emf$.}
	\label{tikz:LCR_circuit}
\end{figure}

Виділимо тепер у циркуляції окремі ділянки, що містять по
одному суттєвому елементу кола:
\begin{equation}\label{eq:circulation_decomposed}
	\oint\limits_{\{12341\}} E_l \, dl
	= \int\limits_{(12)} E_l \, dl
	+ \int\limits_{(23)} E_l \, dl
	+ \int\limits_{(34)} E_l \, dl
	+ \int\limits_{(41)} E_l \, dl.
\end{equation}

Доданок $\int\limits_{(41)} E_l \, dl = 0$, оскільки на цій ділянці немає елементів, що
змінюють потенціал. Доданок
\begin{equation}\label{eq:phi12}
	\int\limits_{(12)} E_l \, dl = \phi_1 - \phi_2
\end{equation}
визначає зміну потенціалу на ділянці 12 і враховує стороннє
джерело (стік) енергії --- електрорушійну силу. ЕРС вводиться
рівністю
\begin{equation}\label{eq:EMF_definition}
	\emf = \phi_2 - \phi_1
\end{equation}
і дає \emph{приріст} потенціалу при переході цього елемента.

Розглянемо ділянку контуру 23. На цій ділянці з не врахованих
раніше елементів присутні тільки провідники з опором $R$, тому
\begin{equation}\label{eq:IR}
	\int\limits_{(23)} E_l \, dl = IR.
\end{equation}

Ділянка контуру 34 містить конденсатор. Вважаємо, що на верхній
пластині є заряд $(+q)$, а на нижній $(-q)$. Тоді струм у додатному
напрямку пов'язаний із зарядом $q$ співвідношенням $I = dq/dt$ і
відповідає заряджанню верхньої пластини конденсатора. За означенням
напруга на конденсаторі є $V = \phi_3 - \phi_4$. З урахуванням цього маємо
\begin{equation}\label{eq:capacitor_voltage}
	\int\limits_{(34)} E_l \, dl = \phi_3 - \phi_4 = V = \frac{q}{C},
\end{equation}
де введено ємність конденсатора $C = q/V$.

З урахуванням сказаного отримуємо
\begin{equation}\label{eq:circuit_equation_raw}
	-\emf + IR + \frac{q}{C} = -\frac{d\Phi}{dt}.
\end{equation}
Перепишемо це рівняння, маючи на увазі зв'язок магнітного потоку і струму
$\Phi = LI$:
\begin{equation}\label{eq:circuit_equation}
	\tcbhighmath{
		L\frac{dI}{dt} + IR + \frac{q}{C} = \emf.
	}
\end{equation}

%% --------------------------------------------------------
\subsection*{Рівняння вільних коливань}
%% --------------------------------------------------------

Нехай у замкненому колі, показаному на рис.~\ref{tikz:LCR_circuit}, відсутнє
джерело ЕРС. Тоді рівняння \eqref{eq:circuit_equation} набирає вигляду
\begin{equation}\label{eq:free_oscillation_raw}
	L\frac{dI}{dt} + IR + \frac{q}{C} = 0.
\end{equation}
Введемо позначення
\begin{equation}\label{eq:definitions}
	2\gamma = \frac{R}{L}, \qquad \omega_0^2 = \frac{1}{LC}
\end{equation}
і врахуємо, що $I = dq/dt$. Тоді отримаємо таке рівняння:
\begin{equation}\label{eq:free_oscillation}
	\tcbhighmath{
		\ddot{q} + 2\gamma\dot{q} + \omega_0^2 q = 0.
	}
\end{equation}
Це рівняння описує поведінку заряду на конденсаторі за
відсутності сторонніх джерел енергії і називається \emphx{рівнянням вільних коливань}\index{Коливання!вільні}.

%% --------------------------------------------------------
\subsection*{Гармонічні коливання}
%% --------------------------------------------------------

Нехай опір кола дорівнює нулю: $R = 0$. Тоді $\gamma = 0$, і ми
приходимо до рівняння гармонічного осцилятора:
\begin{equation}\label{eq:harmonic_equation}
	\ddot{q} + \omega_0^2 q = 0.
\end{equation}
Розв'язок цього рівняння має вигляд
\begin{equation}\label{eq:harmonic_solution}
	q(t) = q_m \cos(\omega_0 t + \phi_0).
\end{equation}
Значення амплітуди $q_m$ і початкової фази $\phi_0$ коливань визначаються
з початкових умов. Частота коливань є
\begin{equation}\label{eq:frequency}
	\omega_0 = \frac{1}{\sqrt{LC}}.
\end{equation}
Відповідно період коливань дорівнює
\begin{equation}\label{eq:Thomson_formula}
	\tcbhighmath{
		T = \frac{2\pi}{\omega_0} = 2\pi\sqrt{LC}.
	}
\end{equation}
Це співвідношення називається \emphx{формулою Томсона}\index{Томсона формула}.

Система, що здійснює коливання, називається \emphx{осцилятором}. Осцилятор називається
\emphx{гармонічним}, якщо коливання, які він здійснює, є гармонічними, тобто описуються
формулою
\begin{equation}\label{eq:harmonic_general}
	q(t) = q_m \cos(\omega t + \phi_0).
\end{equation}
Величина $q_m$ в цьому співвідношенні називається \emphx{амплітудою}, а аргумент
косинуса $\phi = \omega t + \phi_0$ --- \emphx{фазою коливань}.

Рівняння \eqref{eq:harmonic_equation} тотожне за формою рівнянню руху
пружинного маятника без тертя $\ddot{x} + \omega_0^2 x = 0$. Це не
випадковий збіг, а наслідок спільної структури обох систем: у кожній з
них є \enquote{інерційний} елемент, що накопичує кінетичну енергію, і
\enquote{пружний} елемент, що накопичує потенціальну енергію. Нижче наведено
відповідність величин, яка допомагає переносити всі результати,
отримані для контуру, на механічний осцилятор (і навпаки).

\begin{center}
\renewcommand{\arraystretch}{1.4}
	\begin{tblr}{
		colspec={X[l, m]X[c,m]},
        row{odd} = {gray!10},
		row{1} = {c, m, fg=white, bg=themecolorlight, font=\bfseries},
		hlines,
		vlines
		}
\textbf{Коливальний контур} & \textbf{Пружинний маятник} & \textbf{Фізична роль} \\
Заряд $q$ & Зміщення $x$ & Динамічна змінна \\
Струм $I = \dot{q}$ & Швидкість $v = \dot{x}$ & Швидкість зміни \\
Індуктивність $L$ & Маса $m$ & Інертність \\
%Опір $R$ & Тертя $r$ & Дисипація енергії \\
$1/C$ & Жорсткість $k$ & Повертальна дія \\
\end{tblr}
\end{center}

Порівняння станів системи у чотирьох фазах коливань наведено на
рис.~\ref{fig:oscillation_phases}.

\begin{figure}[htbp]\centering
\localinput{oscillations.tikz}
\caption{Фази коливань електричного та магнітного полів у коливальному
контурі та координати пружинного маятника}
\label{fig:oscillation_phases}
\end{figure}

Важливо усвідомити, що в контурі коливається не лише заряд на
конденсаторі. Разом із зарядом періодично змінюються \emph{обидва}
поля: електричне поле в конденсаторі та магнітне поле в котушці. Ці
поля відіграють ту саму роль, що потенціальна та кінетична енергія
маятника. У фазі $0^\circ$ конденсатор заряджений максимально, струм
дорівнює нулю, і вся енергія зосереджена в електричному полі --- це
повний аналог маятника, відхиленого від рівноваги, коли вся енергія
потенціальна. У фазі $90^\circ$ конденсатор розряджений, струм
максимальний, і енергія перейшла в магнітне поле котушки --- аналог
проходження маятником положення рівноваги з максимальною швидкістю,
коли енергія суто кінетична. У фазах $180^\circ$ і $270^\circ$ картина
повторюється з протилежним знаком: конденсатор перезаряджається, а
струм змінює напрямок. Отже, вільні коливання в контурі --- це
періодичне взаємне перетворення енергій електричного та магнітного
полів, точно так само, як коливання маятника --- це періодичне
перетворення потенціальної енергії пружини в кінетичну енергію
тягарця.

%% --------------------------------------------------------
\subsection*{Згасаючі коливання}
%% --------------------------------------------------------

Розглянемо\index{Коливання!згасаючі} загальний випадок рівняння вільних коливань
\begin{equation}\label{eq:damped_equation}
	\ddot{q} + 2\gamma\dot{q} + \omega_0^2 q = 0.
\end{equation}
Введемо комплекснозначну функцію $\hat{q}(t)$, таку, що $\operatorname{Re}\hat{q}(t) = q(t)$, і
яка підпорядковується тому самому рівнянню, що й $q(t)$:
\begin{equation}\label{eq:complex_damped}
	\frac{d^2\hat{q}}{dt^2} + 2\gamma\frac{d\hat{q}}{dt} + \omega_0^2\hat{q} = 0.
\end{equation}
Будемо шукати розв'язок у формі
\begin{equation}\label{eq:complex_ansatz}
	\hat{q}(t) = q_m e^{i\Omega t}.
\end{equation}
Підстановка цього виразу в \eqref{eq:complex_damped} дає
\begin{equation}\label{eq:characteristic}
	\left[-\Omega^2 + 2i\gamma\Omega + \omega_0^2\right]\hat{q} = 0,
\end{equation}
що приводить до квадратного рівняння для $\Omega$:
\begin{equation}\label{eq:Omega_roots}
	\Omega^2 - 2i\gamma\Omega - \omega_0^2 = 0
	\quad\Rightarrow\quad
	\Omega_{1,2} = i\gamma \pm \sqrt{\omega_0^2 - \gamma^2}.
\end{equation}
Відповідно загальний розв'язок диференціального рівняння \eqref{eq:complex_damped}
є лінійною комбінацією частинних розв'язків, що відповідають двом кореням
цього квадратного рівняння:
\begin{equation}\label{eq:general_solution}
	\hat{q}(t) = q_{10} e^{i\Omega_1 t} + q_{20} e^{i\Omega_2 t}.
\end{equation}
Коефіцієнти $q_{10}$ і $q_{20}$ визначаються з початкових умов.

Розв'язок вихідного рівняння \eqref{eq:free_oscillation} є $q(t) = \operatorname{Re}\hat{q}(t)$.

Наведемо частинні випадки.

1) Слабке згасання: $\omega_0 > \gamma$.

Позначимо $\omega = \sqrt{\omega_0^2 - \gamma^2}$. Тоді $\Omega_{1,2} = i\gamma \pm
\omega$, і розв'язок можна переписати у вигляді
\begin{equation}\label{eq:weak_damping_complex}
	\hat{q}(t) = e^{-\gamma t}\left[q_{10} e^{i\omega t} + q_{20} e^{-i\omega t}\right].
\end{equation}
Маючи на увазі формулу Ейлера
\begin{equation}\label{eq:Euler}
	e^{i\phi} = \cos\phi + i\sin\phi
\end{equation}
і відділяючи дійсну частину $q(t) = \operatorname{Re}\hat{q}(t)$, отримаємо
\begin{equation}\label{eq:weak_damping}
	q(t) = e^{-\gamma t}\left(a\cos\omega t + b\sin\omega t\right) = q_m e^{-\gamma t}\cos(\omega t + \phi_0).
\end{equation}
Розв'язок містить дві довільні сталі $(a, b)$ або $(q_m, \phi_0)$.
Ці пари констант пов'язані співвідношеннями
\begin{equation}\label{eq:constants_relation}
	q_m = \sqrt{a^2 + b^2}, \qquad \tg\phi_0 = -\frac{b}{a}.
\end{equation}

2) Сильне згасання: $\omega_0 < \gamma$.

Позначимо $\Gamma = \sqrt{\gamma^2 - \omega_0^2}$. Тоді $\Omega = i(\gamma \pm \Gamma)$, і розв'язок набирає вигляду
\begin{equation}\label{eq:strong_damping}
	q(t) = e^{-\gamma t}\left[q_{10} e^{\Gamma t} + q_{20} e^{-\Gamma t}\right].
\end{equation}

3) Якщо $\omega_0 = \gamma$, то $\Omega_1 = \Omega_2$, і розв'язок вироджується:
\begin{equation}\label{eq:critical_damping}
	q(t) = e^{-\gamma t}\left(C_1 + C_2 t\right)
\end{equation}
($C_1$ і $C_2$ --- довільні сталі).

На рис.~\ref{tikz:damped} показано графік згасаючих коливань. Такі коливання можна представляти як періодичні (гармонічні) коливання зі спадною амплітудою:
\begin{equation}\label{eq:damped_general}
	q(t) = A(t)\cos(\omega t + \phi_0), \qquad A(t) = q_m e^{-\gamma t}.
\end{equation}
На цьому рисунку часова залежність амплітуди згасаючих коливань показана штриховою лінією.

\begin{figure}[h!]
	\centering
	\localinput{damped_oscillations.tikz}
	\caption{Графік згасаючих коливань. Штриховою лінією показано часову залежність амплітуди $A(t) = q_m e^{-\gamma t}$.}
	\label{tikz:damped}
\end{figure}

Величина $\gamma$ показує, з якою швидкістю зменшується амплітуда коливань $A(t) = q_m e^{-\gamma t}$, і називається \emphx{коефіцієнтом згасання}. Для $LCR$-контуру коефіцієнт згасання рівний $\gamma = R/2L$. Обернена величина $\tau = 1/\gamma$ називається \emphx{часом згасання}, і показує, за який
час амплітуда коливань зменшується в $e$ разів.

%% --------------------------------------------------------
\subsection*{Закон збереження енергії}
%% --------------------------------------------------------

Запишемо диференціальне рівняння згасаючих коливань у такій формі:
\begin{equation}\label{eq:damped_equation_energy}
	L\ddot{q} + \frac{q}{C} = -R\dot{q}.
\end{equation}
Помножимо почленно це рівняння на $\dot{q}$ і врахуємо тотожності
\begin{equation}\label{eq:identities}
	\dot{q}\ddot{q} = \frac{d}{dt}\left(\frac{1}{2}\dot{q}^2\right), \qquad
	\dot{q}q = \frac{d}{dt}\left(\frac{1}{2}q^2\right).
\end{equation}
Оскільки $\dot{q} = I$, то отримуємо рівність
\begin{equation}\label{eq:energy_balance}
	\tcbhighmath{
		\frac{d}{dt}\left(\frac{LI^2}{2} + \frac{q^2}{2C}\right) = -RI^2.
	}
\end{equation}
Величина в дужках
\begin{equation}\label{eq:total_energy}
	W = W_L + W_C, \qquad W_L = \frac{LI^2}{2}, \qquad W_C = \frac{q^2}{2C}
\end{equation}
є повна енергія контуру, запасена в котушці індуктивності ($W_L$) і в конденсаторі ($W_C$). Доданок у правій частині рівності ($-N = -RI^2$) описує, згідно із законом Джоуля--Ленца, втрати енергії. Таким чином, отримуємо
\begin{equation}\label{eq:energy_loss}
	\tcbhighmath{
		\frac{dW}{dt} = -N.
	}
\end{equation}

%% --------------------------------------------------------
\subsection*{Перетворення енергії в контурі без втрат}
%% --------------------------------------------------------

Якщо опір контуру дорівнює нулю, то втрати відсутні,
$N = 0$, і енергія зберігається. При цьому заряд на конденсаторі і струм у
колі змінюються за законом
\begin{equation}\label{eq:lossless_oscillations}
	q(t) = q_m \cos(\omega_0 t), \qquad
	I(t) = \dot{q}(t) = -I_m \sin(\omega_0 t) = I_m \cos\left(\omega_0 t + \frac{\pi}{2}\right),
\end{equation}
де введена амплітуда струму $I_m = \omega_0 q_m$. Звідси видно, що фази $q(t)$ і
$I(t)$ зсунуті на $\pi/2$. Відповідно для доданків в енергії знаходимо:
\begin{equation}\label{eq:WL_WC_oscillations}
	W_L = \frac{LI^2}{2} = \frac{LI_m^2}{2} \cos^2\left(\omega_0 t + \frac{\pi}{2}\right) = \frac{LI_m^2}{2} \frac{1 - \cos(2\omega_0 t)}{2},
\end{equation}
\begin{equation}\label{eq:WC_oscillations}
	W_C = \frac{q^2}{2C} = \frac{q_m^2}{2C} \cos^2(\omega_0 t) = \frac{q_m^2}{2C} \frac{1 + \cos(2\omega_0 t)}{2}.
\end{equation}
Очевидно, що
\begin{equation}\label{eq:energy_conservation}
	\tcbhighmath{
		W_L + W_C = \frac{LI^2}{2} + \frac{q^2}{2C} = \frac{LI_m^2}{2} + \frac{q_m^2}{2C} = \const.
	}
\end{equation}
Крім того, в силу співвідношень $I_m = \omega_0 q_m$, $\omega_0^2 = 1/LC$ отримуємо
\begin{equation}\label{eq:energy_amplitudes}
	W_L^{(0)} = \frac{LI_m^2}{2} = \frac{L\omega_0^2}{2} q_m^2 = \frac{q_m^2}{2C} = W_C^{(0)}.
\end{equation}
Таким чином, запаси енергії в котушці індуктивності і в конденсаторі здійснюють коливання з однаковою амплітудою, але в протифазі,
причому періодично відбувається повне перекачування енергії з індуктивності в ємність і назад.

Якщо контур неідеальний ($R \neq 0$), то енергія втрачається. Максимальна швидкість втрат досягається в ті моменти, коли струм у колі
максимальний, оскільки в ці моменти максимальні джоулеві втрати
($N = RI^2$).

%% --------------------------------------------------------
\subsection*{Характеристики згасаючих коливань}
%% --------------------------------------------------------

Для кількісного опису згасаючих коливань використовують декілька взаємопов'язаних характеристик, кожна з яких описує згасання зі свого боку: через зменшення амплітуди, через втрати енергії, через час життя коливання тощо. Розглянемо основні з них.

\paragraph*{Логарифмічний декремент згасання.}
Розглянемо згасаючі коливання (рис.~\ref{tikz:damped}):
\begin{equation}\label{eq:damped_general_2}
	q(t) = A(t)\cos(\omega t + \phi_0), \qquad A(t) = q_m e^{-\gamma t}.
\end{equation}
Виберемо два послідовних максимуми: $q_n$, $q_{n+1}$. Величина
\begin{equation}\label{eq:log_decrement}
	\tcbhighmath{
		\delta = \ln(q_n / q_{n+1})
	}
\end{equation}
називається \emphx{логарифмічним декрементом}\index{Логарифмічний декремент згасання}. Знайдемо явний вираз
для $\delta$.

\begin{figure}[h!]
	\centering
\localinput{log_decrement.tikz}
	\caption{Згасаючі коливання. До означення логарифмічного декремента.}
	\label{tikz:log_decrement}
\end{figure}

Максимуми функції $q(t)$ визначаються з умови $\dot{q} = 0$. З виразу \eqref{eq:damped_general_2} знаходимо
\begin{equation}\label{eq:derivative_maxima}
	\begin{aligned}
		\dot{q} & = q_m e^{-\gamma t}\left[-\gamma \cos(\omega t + \phi_0) - \omega \sin(\omega t + \phi_0)\right] = \\
		        & = -q_m \omega_0 e^{-\gamma t} \cos(\omega t + \phi_0 + \phi_1).
	\end{aligned}
\end{equation}
Додатковий доданок $\phi_1$ у фазі такий, що
\begin{equation}\label{eq:varphi1}
	\cos\phi_1 = \frac{\gamma}{\sqrt{\gamma^2 + \omega^2}} = \frac{\gamma}{\omega_0},
\end{equation}
\begin{equation}\label{eq:varphi1_sin}
	\sin\phi_1 = -\frac{\omega}{\sqrt{\gamma^2 + \omega^2}} = -\frac{\omega}{\omega_0}.
\end{equation}
Враховано також, що відповідно до означення величини $\omega$ виявляється $\sqrt{\gamma^2 + \omega^2} = \omega_0$.

Покладемо $\dot{q} = 0$. Як видно з \eqref{eq:derivative_maxima}, проміжок часу між
двома сусідніми максимумами становить
\begin{equation}\label{eq:period}
	T = 2\pi/\omega
\end{equation}
і дорівнює періоду коливань множника $\cos(\omega t + \phi_0)$ у виразі для
$q(t)$. Це означає, що згідно з \eqref{eq:damped_general_2}
\begin{equation}\label{eq:maxima_ratio}
	\frac{q_n}{q_{n+1}} = \frac{A(t_n)}{A(t_{n+1})} = e^{\gamma T},
\end{equation}
звідки випливає вираз для логарифмічного декремента:

\begin{equation}\label{eq:delta_explicit}
	\delta = \gamma T = 2\pi\gamma/\omega.
\end{equation}
У випадку слабкого згасання ($\gamma \ll \omega_0$) для $LCR$-контуру знаходимо:
\begin{equation}\label{eq:delta_weak}
	\omega \approx \omega_0 = \frac{1}{\sqrt{LC}}, \qquad
	\gamma = \frac{R}{2L}
	\quad\Rightarrow\quad
	\delta = \frac{2\pi\gamma}{\omega} = \pi R \sqrt{\frac{C}{L}}.
\end{equation}
Смисл логарифмічного декремента такий. За характерний
час згасання коливань $\tau = 1/\gamma$ осцилятор здійснить
\begin{equation}\label{eq:number_oscillations}
	n = \frac{\tau}{T} = \frac{1}{\gamma T} = \frac{1}{\delta}
\end{equation}
коливань. Отже, $\delta = 1/n$.

\paragraph*{Добротність коливальної системи.}
Величина $Q = \pi n$, або
\begin{equation}\label{eq:quality_factor}
	\tcbhighmath{
		Q = \omega/2\gamma
	}
\end{equation}
називається \emphx{добротністю}\index{Добротність} коливальної системи. Очевидний зв'язок
добротності і логарифмічного декремента: $Q = \pi/\delta$. Для послідовного $LCR$-контуру зі слабким згасанням маємо
\begin{equation}\label{eq:Q_LCR}
	\tcbhighmath{
		Q = \frac{1}{R}\sqrt{\frac{L}{C}}.
	}
\end{equation}

За період коливань амплітуда коливань заряду і струму зменшується в
$e^{\gamma T}$ разів. При цьому згідно з \eqref{eq:total_energy} енергія системи зменшується як квадрат амплітуди, тобто в $e^{2\gamma T}$ разів. Це означає, що якщо на початку якогось періоду коливань в контурі запасена енергія $W(t)$, то до початку
наступного циклу в системі залишається енергія $W(t + T) = e^{-2\gamma T}W(t)$. Втрати енергії за цикл становлять
\begin{equation}\label{eq:energy_loss_cycle}
	\Delta W = W(t) - W(t + T) = \left(1 - e^{-2\gamma T}\right)W(t).
\end{equation}
Вважаючи згасання слабким: $2\gamma T \ll 1$, отримуємо $\Delta W = 2\gamma T \cdot W(t)$. Звідси
випливає
\begin{equation}\label{eq:quality_energy}
	\tcbhighmath{
		\frac{W}{\Delta W} = \frac{1}{2\gamma T} = \frac{1}{2\pi}\frac{\omega}{2\gamma} = \frac{1}{2\pi}Q,
	}
\end{equation}
що й визначає енергетичний смисл добротності $Q$. Останню
формулу зазвичай записують у вигляді
\begin{equation}\label{eq:Q_definition}
	\tcbhighmath{
		Q = 2\pi\frac{W}{\Delta W}.
	}
\end{equation}
Цю формулу іноді приймають як означення добротності.

\begin{Attention}
	Добротність --- це, по суті, спосіб виміряти, наскільки \enquote{дзвінкою} є коливальна система: скільки коливань вона встигне зробити, перш ніж розтратити запасену енергію. Величини $Q$ для різних систем відрізняються на порядки. У звичайному $LC$-контурі з дроту й повітряного конденсатора добротність --- це, як правило, кілька десятків чи сотень. У кварцового резонатора наручного годинника вона сягає $10^{4}$--$10^{6}$, тому такий годинник і не поспішає та не запізнюється роками. А рекордсменами є надпровідні резонатори, які використовують у прискорювачах заряджених частинок: їхня добротність доходить до $10^{9}$--$10^{11}$ --- електромагнітне поле в такому резонаторі \enquote{дзвенить}, практично не втрачаючи енергії, мільярди періодів підряд.
\end{Attention}

%% --------------------------------------------------------
\section{Вимушені коливання під дією гармонічної ЕРС}\hypertarget{forced-oscillations}{}
%% --------------------------------------------------------

%% --------------------------------------------------------
\subsection*{Рівняння вимушених коливань}
%% --------------------------------------------------------

Запишемо\index{Коливання!вимушені} рівняння \eqref{eq:circuit_equation} для послідовного $LCR$-контуру за наявності ЕРС $\emf$,
що змінюється за гармонічним законом:
\begin{equation}\label{eq:forced_equation}
	L\ddot{q} + R\dot{q} + \frac{1}{C}q = \emf(t), \qquad
	\emf(t) = \emf_m \cos\omega t.
\end{equation}
Переходячи від заряду на конденсаторі до напруги ($V = q/C$) і використовуючи введені раніше позначення \eqref{eq:definitions} ($2\gamma = R/L$, $\omega_0^2 = 1/LC$), перепишемо це рівняння у вигляді
\begin{equation}\label{eq:forced_equation_V}
	\ddot{V} + 2\gamma\dot{V} + \omega_0^2 V = \omega_0^2 \emf_m \cos\omega t.
\end{equation}

\begin{Attention}
	Мати справу з гармонічними функціями не так просто; набагато зручніше працювати з експоненційними, які легко диференціювати й інтегрувати. Ключем, що пов'язує їх, є формула Ейлера
	\begin{equation*}
		e^{i\phi} = \cos\phi + i\sin\phi.
	\end{equation*}

	Нехай дійсна величина $x(t)$ здійснює гармонічні коливання:
	\begin{equation}\label{eq:harmonic_real}
		x(t) = x_0 \cos(\omega t + \phi_0).
	\end{equation}
	Введемо комплексну функцію
	\begin{equation}\label{eq:complex_function}
		z(t) = z_0 e^{i\omega t},
	\end{equation}
	таку що $x(t) = \operatorname{Re} z(t)$. Коефіцієнт $z_0$ називається \emphx{комплексною амплітудою} величини $x$. У розглянутому прикладі
	\begin{equation}\label{eq:complex_amplitude}
		z_0 = x_0 e^{i\phi_0}.
	\end{equation}
	Отже, комплексна амплітуда несе інформацію як про амплітуду $x_0$, так і про початкову фазу $\phi_0$ коливання.

	Нагадаємо, що будь-яке комплексне число $z = a + ib$ можна записати як $z = |z| e^{i\phi}$, де $|z| = \sqrt{a^2 + b^2}$ --- модуль, а $\phi$ --- фаза, що визначається як $\tg\phi = b/a$; при цьому $\operatorname{Re} z = |z|\cos\phi$. Саме ці співвідношення й лежать в основі комплексного представлення: амплітуда $x_0$ відіграє роль модуля, а початкова фаза $\phi_0$ --- роль фази комплексної амплітуди $z_0 = x_0 e^{i\phi_0}$.
\end{Attention}

Далі, для розв'язання отриманого рівняння перейдемо до комплексного
представлення коливань, поклавши
\begin{equation}\label{eq:complex_representation}
	V(t) \to \hat{V}(t) = \operatorname{Re}\hat{V}(t) + i\operatorname{Im}\hat{V}(t), \qquad
	V(t) = \operatorname{Re}\hat{V}(t);
\end{equation}
\begin{equation}\label{eq:complex_EMF}
	\emf(t) \to \hat{\emf}(t) = \emf_m e^{i\omega t}, \qquad
	\emf(t) = \operatorname{Re}\hat{\emf}(t).
\end{equation}
Тоді замість \eqref{eq:forced_equation_V} будемо мати справу з рівнянням
\begin{equation}\label{eq:complex_forced}
	\ddot{\hat{V}} + 2\gamma\dot{\hat{V}} + \omega_0^2 \hat{V} = \omega_0^2\emf_m e^{i\omega t}.
\end{equation}
Розв'язок останнього рівняння будемо шукати у вигляді
\begin{equation}\label{eq:solution_decomposition}
	\hat{V}(t) = \hat{V}_s(t) + \hat{V}_f(t),
\end{equation}
де доданок $\hat{V}_s(t)$ описує вільні коливання (\enquote{$s$} = \emph{self}), а доданок $\hat{V}_f(t)$ --- вимушені коливання (\enquote{$f$} = \emph{forced}):
\begin{equation}\label{eq:forced_system}
	\begin{cases}
		\ddot{\hat{V}}_s + 2\gamma\dot{\hat{V}}_s + \omega_0^2 \hat{V}_s = 0, \\[6pt]
		\ddot{\hat{V}}_f + 2\gamma\dot{\hat{V}}_f + \omega_0^2 \hat{V}_f = \omega_0^2\emf_m e^{i\omega t}.
	\end{cases}
\end{equation}

Якщо коефіцієнт згасання ненульовий ($\gamma \neq 0$), то вільні коливання з часом згаснуть: $V_s(t) \to 0$, і залишаться тільки вимушені
коливання. Знайдемо розв'язок $\hat{V}_f(t)$, що відповідає вимушеним коливанням.

Будемо шукати розв'язок у вигляді
\begin{equation}\label{eq:ansatz_forced}
	\hat{V}_f(t) = \hat{V}_m e^{i\Omega t}.
\end{equation}
Підстановка цього виразу в друге рівняння в \eqref{eq:forced_system} дає
\begin{equation}\label{eq:substitution_forced}
	\left(\omega_0^2 - \Omega^2 + 2i\gamma\Omega\right)\hat{V}_m e^{i\Omega t} = \omega_0^2\emf_m e^{i\omega t}.
\end{equation}
Оскільки така рівність повинна виконуватися в будь-який момент часу, то звідси випливає $\Omega = \omega$ і
\begin{equation}\label{eq:V0_frequency_response}
	\tcbhighmath{
		\hat{V}_m = \frac{\omega_0^2}{\omega_0^2 - \omega^2 + 2i\gamma\omega}\emf_m.
	}
\end{equation}
Отримана залежність називається \emphx{частотною характеристикою} контуру.

Покладемо
\begin{equation}\label{eq:polar_form}
	\omega_0^2 - \omega^2 + 2i\gamma\omega = \rho e^{-i\phi_0}.
\end{equation}
Звідси знаходимо
\begin{equation}\label{eq:rho_phi}
	\rho = \sqrt{\left(\omega_0^2 - \omega^2\right)^2 + 4\gamma^2\omega^2},
\end{equation}
\begin{equation}\label{eq:tg_phi0}
	\tg\phi_0 = \frac{2\gamma\omega}{\omega^2 - \omega_0^2}.
\end{equation}
З урахуванням цього розв'язок, що описує вимушені коливання напруги, набирає вигляду
\begin{equation}\label{eq:forced_voltage}
	\tcbhighmath{
		V_f = V_m \cos(\omega t + \phi_0), \qquad
		V_m = \frac{\omega_0^2}{\sqrt{\left(\omega_0^2 - \omega^2\right)^2 + 4\gamma^2\omega^2}}\emf_m.
	}
\end{equation}

\begin{Historical}
	Перехід від диференціального рівняння\eqref{eq:forced_equation} до звичайної алгебри з комплексними числами \eqref{eq:substitution_forced}, яким ми щойно скористалися, здавався інженерам XIX століття далеко не очевидним. До 1890-х років розрахунок кожного нового кола змінного струму означав щоразу розв'язувати диференціальні рівняння вручну, що на практиці ставало непідйомним для розгалужених промислових мереж. Усе змінив інженер компанії General Electric Чарльз Протеус Штейнмец: у доповіді \enquote{Complex Quantities and Their Use in Electrical Engineering} (1893) він показав, що достатньо замінити напруги й струми комплексними амплітудами --- і розрахунок кола перетворюється на просту алгебру, якою ми користуємося у цьому розділі. Метод Штейнмеца настільки прискорив і спростив проєктування електричних мереж, що саме йому значною мірою завдячують швидким розгортанням систем передачі й розподілу електроенергії на змінному струмі наприкінці XIX --- на початку XX століття.
\end{Historical}

%% --------------------------------------------------------
\subsection*{Амплітудно-частотна характеристика}
%% --------------------------------------------------------

На рис.~\ref{tikz:amplitude_response} показана залежність амплітуди вимушених коливань від частоти $\omega$ зовнішньої ЕРС, яка називається \emphx{амплітудно-частотною}\index{Амплітудно-частотна характеристика} (або \emphx{амплітудною}) характеристикою.

\begin{figure}[h!]
	\centering
	\localinput{amplitude_response.tikz}
	\caption{Амплітудно-частотна характеристика контуру. Залежність амплітуди вимушених коливань $V_m$ від частоти $\omega$ зовнішньої ЕРС.}
	\label{tikz:amplitude_response}
\end{figure}

Як видно з рис.~\ref{tikz:amplitude_response}, при не надто великому згасанні спостерігається сильне зростання амплітуди при наближенні частоти
зовнішньої сили (ЕРС) до деякої характерної частоти. Це явище
називається \emphx{резонансом}\index{Резонанс}, а частота $\omega_m$, при якій амплітуда досягає
максимуму, --- \emphx{резонансною частотою}. Докладніше про резонанс ітиметься нижче, а також у \S~\ref{sec:ac_resonance}, присвяченому змінному струму, де ми побачимо його з боку імпедансів.

Якщо згасання відсутнє: $\gamma = 0$, то резонансна частота збігається з власною частотою осцилятора: $\omega_m = \omega_0$. При цьому
\begin{equation}\label{eq:resonance_no_damping}
	V_m = \frac{\omega_0^2}{\left|\omega_0^2 - \omega^2\right|}\emf_m.
\end{equation}
Якщо ж $\gamma \neq 0$, то резонансну частоту можна знайти з \eqref{eq:forced_voltage}, поклавши
$dV_m/d\omega = 0$:
\begin{equation}\label{eq:resonance_frequency}
	\tcbhighmath{
		\omega_m = \sqrt{\omega_0^2 - 2\gamma^2}.
	}
\end{equation}
Якщо $\gamma \geqslant \omega_0/\sqrt{2}$, то з ростом частоти ЕРС амплітуда коливань монотонно зменшується при всіх частотах.

%% --------------------------------------------------------
\subsection*{Фазово-частотна характеристика}
%% --------------------------------------------------------

Як видно з \eqref{eq:rho_phi}, \eqref{eq:forced_voltage}, коливання напруги $V_f(t)$ зсунуті по фазі відносно коливань ЕРС $\emf(t)$. На рис.~\ref{tikz:phase_response} показана залежність фазового зсуву $\phi_0$ від частоти ЕРС, яка називається \emphx{фазово-частотною}\index{Фазово-частотна характеристика} (або \emphx{фазовою}) характеристикою.

\begin{figure}[h!]
	\centering
	\localinput{phase_response.tikz}
	\caption{Фазово-частотна характеристика контуру. Залежність фазового зсуву $\phi_0$ від частоти $\omega$ зовнішньої ЕРС.}
	\label{tikz:phase_response}
\end{figure}

%% --------------------------------------------------------
\subsection*{Резонанс}
%% --------------------------------------------------------

У багатьох практично важливих випадках згасання в коливальних системах мале, $\gamma \ll \omega_0$. Знайдемо для цього випадку характеристики
резонансної кривої (рис.~\ref{tikz:resonance_curve}).

Згідно з \eqref{eq:forced_voltage} резонансна частота $\omega_m \approx \omega_0$. Максимум амплітуди досягається при $\omega \approx \omega_0$ і становить
\begin{equation}\label{eq:Vmax}
	\tcbhighmath{
		V_{\max} \approx \frac{\omega_0}{2\gamma}\emf_m.
	}
\end{equation}

При $\omega \to 0$ амплітуда $V_m$ прямує до границі
\begin{equation}\label{eq:static_limit}
	V_{\text{ст}} = \emf_m.
\end{equation}
Цей результат відображає відгук на статичну дію
$\emf = \emf_m = \const$ і при $\dot{V} = 0$, $\ddot{V} = 0$ прямо випливає з \eqref{eq:forced_equation_V}.

\begin{figure}[h!]
	\centering
	\localinput{resonance_curve.tikz}
	\caption{Характеристики резонансної кривої: $V_{\max}$ --- максимум амплітуди, $V_{\text{ст}}$ --- статична амплітуда, $\Delta\omega$ --- ширина резонансу.}
	\label{tikz:resonance_curve}
\end{figure}

Знайдемо відношення $V_{\max}/V_{\text{ст}}$. Згідно з \eqref{eq:Vmax}, \eqref{eq:static_limit} маємо
\begin{equation}\label{eq:Q_resonance}
	\tcbhighmath{
		\frac{V_{\max}}{V_{\text{ст}}} = \frac{\omega_0}{2\gamma} = Q.
	}
\end{equation}
Це відношення збігається з добротністю $Q$ розглядуваної системи.
Таким чином, в резонансі амплітуда коливань в $Q$ разів більша, ніж
статичне відхилення: $V_{\max} = QV_{\text{ст}}$.

Амплітуда вимушених коливань зменшується в міру віддалення від
резонансної частоти. Знайдемо ширину резонансу $\Delta\omega = \omega_2 - \omega_1$, де $\omega_1$ і
$\omega_2$ --- частоти нижче і вище резонансної, при яких вона зменшується в
$\sqrt{2}$ разів:
\begin{equation}\label{eq:half_power}
	V_m(\omega) = V_{\max}/\sqrt{2}.
\end{equation}

Будемо вважати згасання слабким. Оскільки при цьому $\omega_m \approx \omega_0$,
$V_{\max} \approx (\omega_0/2\gamma)\emf_m$, то згідно з \eqref{eq:forced_voltage} для знаходження шуканих частот маємо рівняння
\begin{equation}\label{eq:half_power_equation}
	\frac{\emf_m}{\sqrt{\left(\omega_0^2 - \omega^2\right)^2 + 4\gamma^2\omega^2}}
	= \frac{1}{\sqrt{2}}\frac{\emf_m}{2\gamma\omega_0},
\end{equation}
або
\begin{equation}\label{eq:half_power_squared}
	\begin{aligned}
		\left(\omega_0^2 - \omega^2\right)^2 + 4\gamma^2\omega^2 & = 8\gamma^2\omega_0^2 \quad\Rightarrow                                      \\
		\Rightarrow\quad \left(\omega_0^2 - \omega^2\right)^2    & = 4\gamma^2\left(2\omega_0^2 - \omega^2\right) \approx 4\gamma^2\omega_0^2.
	\end{aligned}
\end{equation}
Тут враховано, що в правій частині цієї рівності при малих $\gamma$ можна
покласти $\omega \approx \omega_0$. Таким чином, отримуємо
\begin{equation}\label{eq:frequency_shift}
	\omega_0^2 - \omega^2 \approx \pm 2\gamma\omega_0
	\quad\Rightarrow\quad
	\omega^2 = \omega_0^2\left(1 \mp \frac{2\gamma}{\omega_0}\right),
\end{equation}
звідки випливає:
\begin{equation}\label{eq:omega1_omega2}
	\omega_1 = \omega_0 - \gamma, \qquad \omega_2 = \omega_0 + \gamma.
\end{equation}
Отже, ширина резонансної кривої по рівню $1/\sqrt{2}$ становить
\begin{equation}\label{eq:resonance_width}
	\tcbhighmath{
		\Delta\omega = \omega_2 - \omega_1 = 2\gamma.
	}
\end{equation}
Цей результат можна переписати в іншому вигляді, використовуючи поняття
добротності:
\begin{equation}\label{eq:resonance_width_Q}
	\tcbhighmath{
		\Delta\omega = \frac{\omega_0}{Q}.
	}
\end{equation}
Отже, чим вище добротність системи, тим вужча резонансна
крива.

Тут ми розглянули резонанс у послідовному $LCR$-контурі як явище вимушених коливань --- через амплітуду напруги на конденсаторі. У \S~\ref{sec:ac_resonance}, присвяченому змінному струму, ми побачимо, як той самий результат випливає з аналізу імпедансів, і поширимо його на довільні розгалужені кола.

%% --------------------------------------------------------
\section{Змінний струм}\label{sec:ac}\hypertarget{ac-current}{}
%% --------------------------------------------------------

Змінним струмом називають струм, миттєве значення якого періодично змінюється з часом. Найважливіший випадок --- гармонічний змінний струм, який виникає в колі під дією гармонічної ЕРС. По суті, це ті самі вимушені коливання, які ми розглянули в попередньому розділі. Там ми вже ввели комплексні амплітуди як зручний спосіб опису гармонічних коливань. Тепер ми застосуємо цей метод до розрахунку довільних кіл змінного струму --- як нерозгалужених, так і розгалужених. Виявиться, що всі правила --- закон Ома, правила Кірхгофа, складання опорів --- зберігають свою форму, якщо замінити активний опір $R$ на комплексний імпеданс $Z$.


%% --------------------------------------------------------
\subsection*{Закон Ома в комплексній формі. Імпеданс}
%% --------------------------------------------------------

У попередньому розділі ми знайшли комплексну амплітуду напруги на конденсаторі~\eqref{eq:V0_frequency_response}. Тепер нас цікавить не напруга, а струм у колі. Оскільки $I = \dot{q}$ і $V = q/C$, комплексні амплітуди цих величин пов'язані співвідношенням
\begin{equation}\label{eq:I_from_V}
	\hat{I}_m = i\omega C \hat{V}_m.
\end{equation}
Підставляючи сюди \eqref{eq:V0_frequency_response} і враховуючи, що $\omega_0^2 = 1/(LC)$, отримуємо
\begin{equation}\label{eq:I0_explicit}
	\hat{I}_m = \frac{i\omega / L}{\omega_0^2 - \omega^2 + 2i\gamma\omega}\emf_m.
\end{equation}
Перепишемо знаменник у вигляді
\begin{equation}\label{eq:denominator_impedance}
	\omega_0^2 - \omega^2 + 2i\gamma\omega = \frac{i\omega}{L}\left[R + i\left(\omega L - \frac{1}{\omega C}\right)\right],
\end{equation}
що легко перевірити розкриттям дужок. Тоді
\begin{equation}\label{eq:Ohm_complex_derived}
	\tcbhighmath{
		\hat{I}_m = \frac{\emf_m}{R + i\left(\omega L - \dfrac{1}{\omega C}\right)}.
	}
\end{equation}
Отже, комплексна амплітуда струму пов'язана з комплексною амплітудою ЕРС тим самим співвідношенням, що й у законі Ома для постійного струму, --- але замість активного опору $R$ тут стоїть величина
\begin{equation}\label{eq:impedance}
	\tcbhighmath{
		Z = R + i\left(\omega L - \frac{1}{\omega C}\right),
	}
\end{equation}
яка називається \emphx{імпедансом}\index{Імпеданс}\footnote{Імпеданс --- від англ. \emph{impede} --- заважати, перешкоджати.}, або \emphx{комплексним опором}, електричного кола. У компактному вигляді отримане співвідношення записують як
\begin{equation}\label{eq:Ohm_complex}
	\tcbhighmath{
		\hat{I}_m Z = \emf_m,
	}
\end{equation}
і називають \emphx{законом Ома в комплексній формі}\index{Ома закон!у комплексній формі}: комплексна амплітуда струму дорівнює комплексній амплітуді ЕРС, поділеній на імпеданс. Це той самий закон Ома, що й для постійного струму, лише записаний для комплексних амплітуд гармонічних величин.

%\begin{figure}[h!]
%	\centering
%	\begin{circuitikz}[european resistors, scale=1.5, circuitikz/bipoles/length=2.1cm, thick]
%		\coordinate (A) at (0,0);
%		\coordinate (B) at (4,0);
%		\coordinate (C) at (4,3);
%		\coordinate (D) at (0,3);
%
%		\coordinate (T1) at ($(A)!0.4!(B)$);
%		\coordinate (T2) at ($(A)!0.6!(B)$);
%
%		\draw (A) to[short, -o] (T1) node[above] {};
%		\draw (B) to[short, -o] (T2) node[above] {};
%		\path (T1) -- (T2) node[midway, below=2pt] {$\emf$};
%
%		\draw (B) to[L, mirror,  l_=$L$, color=blue] (C);
%		\draw[fill=green!50!black] (C) to[R, l_=$R$, color=green!50!black] (D);
%
%		\draw[->, >=latex] ($(C)!0.325!(D) + (0,-0.5)$) -- ($(C)!0.675!(D) + (0,-0.5)$)
%		node[midway, below] {$I$};
%
%		\draw (D) to[C, l_=$C$, color=red] (A);
%
%		\node at ($(D)!0.3!(A) + (-0.4,0)$) {$+q$};
%		\node at ($(D)!0.7!(A) + (-0.4,0)$) {$-q$};
%	\end{circuitikz}
%	\caption{Контур, що включає конденсатор, індуктивність, опір і стороннє джерело ЕРС. Стрілка поряд зі струмом вказує додатний напрямок обходу контуру.}
%	\label{tikz:LCR_circuit}
%\end{figure}

Величину
\begin{equation}\label{eq:active_resistance}
	\operatorname{Re} Z = R
\end{equation}
називають \emphx{активним} (або \emphx{омічним}) \emphx{опором}, а величину
\begin{equation}\label{eq:reactive_resistance}
	\operatorname{Im} Z = \omega L - \frac{1}{\omega C}
\end{equation}
--- \emphx{реактивним опором}\index{Опір електричний!реактивний}. Активний опір приводить до втрат енергії в~контурі, тоді як реактивний до втрат не призводить: енергія лише періодично перекачується між джерелом і~полями котушки та конденсатора.

Імпеданс $Z$ описує коло в цілому. Кожен з елементів кола має свій власний імпеданс:
\begin{equation}\label{eq:element_impedances}
	Z_R = R, \qquad Z_C = \frac{1}{i\omega C} = -\frac{i}{\omega C}, \qquad Z_L = i\omega L.
\end{equation}
Звідси видно, що імпеданс резистора --- дійсна величина, а імпеданси конденсатора й котушки --- уявні.  Повний імпеданс послідовного $LCR$-контуру \eqref{eq:impedance} є сумою імпедансів окремих елементів:
\begin{equation}\label{eq:impedance_sum}
	Z = Z_R + Z_C + Z_L.
\end{equation}
Це співвідношення за формою повністю збігається з правилом складання опорів для послідовного кола постійного струму: там теж $R = R_1 + R_2 + \dots$. Різниця лише в тому, що тепер опори --- комплексні величини, і замість простого додавання дійсних чисел доводиться додавати комплексні, враховуючи їхні фази. Саме в цьому й полягає основна зручність методу комплексних амплітуд: усі правила для кіл постійного струму зберігають свою форму, якщо замінити $R$ на $Z$.

Запишемо струм у вигляді $I = I_m e^{i\omega t}$ і помножимо його на імпеданс кожного елемента:
\begin{align}
	U_R &= R I_m e^{i\omega t}, \label{eq:UR_phase} \\[4pt]
	U_C &= -\frac{i}{\omega C} I_m e^{i\omega t}
	      = \frac{I_m}{\omega C} e^{i(\omega t - \pi/2)}, \label{eq:UC_phase} \\[4pt]
	U_L &= i\omega L I_m e^{i\omega t}
	      = \omega L I_m e^{i(\omega t + \pi/2)}. \label{eq:UL_phase}
\end{align}
Звідси видно, що напруга на резисторі збігається за фазою зі струмом, на конденсаторі --- відстає на $\pi/2$, а на котушці --- випереджає на $\pi/2$. Множення на $i$ у комплексній площині означає поворот на $+\pi/2$, а на $-i$ --- на $-\pi/2$; саме тому реактивні елементи зсувають фазу, а активний --- ні.

\subsection*{Векторні діаграми}

Комплексні амплітуди можна зображати векторами на комплексній площині. Тоді додавання струмів і напруг стає додаванням векторів, а закон Ома \eqref{eq:Ohm_complex} --- твердженням про їхню \emph{векторну} суму. Така побудова називається \emphx{векторною діаграмою}\index{Векторна діаграма}.

На рис.~\ref{tikz:vector_diagram_full} показано векторну діаграму для послідовного $LCR$-контуру. Вектор струму $I$ напрямлений уздовж дійсної осі; доданки $RI$, $i\omega L I$ і $-iI/(\omega C)$ утворюють прямокутний трикутник, гіпотенуза якого дорівнює вектору ЕРС $\emf$. З діаграми відразу видно, що
\[
	\emf_m^2 = R^2 I_m^2 + \left(\omega L - \frac{1}{\omega C}\right)^2 I_m^2,
\]
--- тобто векторна діаграма дає той самий результат, що й алгебраїчний підхід, але без обчислень. Саме тому нею зручно користуватися для швидкої якісної оцінки: вона показує, як співвідносяться між собою напруги на елементах і як зсунуті їхні фази.

\begin{figure}[h!]
	\centering
	\localinput{vector_diagram_full.tikz}
	\caption{Векторна діаграма для $LCR$-контуру, показаного на рис.~\ref{tikz:LCR_circuit}, для випадку $\omega L > 1/\omega C$.}
	\label{tikz:vector_diagram_full}
\end{figure}

%% --------------------------------------------------------
\subsection*{Потужність змінного струму}
%% --------------------------------------------------------

Відповідно до закону Джоуля--Ленца (в інтегральній формі)
потужність, що розсіюється струмом у провідному середовищі, рівна $P = \emf I$. У
випадку змінних періодично змінюваних струмів і напруг
інтерес представляє середнє значення потужності за період коливань:
\begin{equation}\label{eq:average_power}
	\overline{P} = \overline{\emf I} = \frac{1}{T}\int\limits_{t}^{t+T} \emf(t_1)I(t_1) \, dt_1.
\end{equation}
У цю формулу входять дійсні величини. Тому користуватися
комплексним представленням слід більш обережно. Покладемо
\begin{equation}\label{eq:real_parts}
	\emf(t) = \frac{1}{2}\left(\hat{\emf}(t) + \hat{\emf}^*(t)\right), \qquad
	I(t) = \frac{1}{2}\left(\hat{I}(t) + \hat{I}^*(t)\right).
\end{equation}
Звідси знаходимо
\begin{equation}\label{eq:product_expansion}
	\begin{aligned}
		\emf(t)I(t) & = \frac{1}{4}\left(\hat{\emf}(t) + \hat{\emf}^*(t)\right)\left(\hat{I}(t) + \hat{I}^*(t)\right) =                                         \\
		            & = \frac{1}{4}\left[\hat{\emf}(t)\hat{I}(t) + \hat{\emf}^*(t)\hat{I}^*(t) + \hat{\emf}(t)\hat{I}^*(t) + \hat{\emf}^*(t)\hat{I}(t)\right] = \\
		            & = \frac{1}{2}\operatorname{Re}\left[\hat{\emf}(t)\hat{I}(t) + \hat{\emf}(t)\hat{I}^*(t)\right].
	\end{aligned}
\end{equation}
Для знаходження середнього врахуємо, що
\begin{equation}\label{eq:complex_amplitudes_time}
	\hat{\emf}(t) = \emf_m e^{i\omega t}, \qquad
	\hat{I}(t) = I_m e^{i\omega t} e^{i\phi_0}
\end{equation}
(амплітуди $\emf_m$ і $I_m$ тут вважаються дійсними величинами). Це
означає, що доданок $\hat{\emf}(t)\hat{I}(t)$ осцилює з частотою $2\omega$ і його середнє
значення за період $T = 2\pi/\omega$ рівне нулю:
\begin{equation}\label{eq:fast_oscillating_term}
	\overline{\hat{\emf}(t)\hat{I}(t)} = \emf_m I_m e^{i\phi_0} \overline{\exp[2i\omega t]}
	= \emf_m I_m e^{i\phi_0} \frac{1}{T}\int\limits_{t}^{t+T} \exp[2i\omega t_1] \, dt_1 = 0.
\end{equation}
Аналогічно встановлюється, що
\begin{equation}\label{eq:conjugate_term}
	\overline{\hat{\emf}^*(t)\hat{I}^*(t)} = 0.
\end{equation}
В результаті з \eqref{eq:average_power} отримуємо
\begin{equation}\label{eq:average_power_complex}
	\tcbhighmath{
		\overline{P} = \frac{1}{2}\operatorname{Re}\left[\hat{\emf}(t)\hat{I}^*(t)\right].
	}
\end{equation}
З урахуванням зсуву фаз між струмом і ЕРС звідси знаходимо
\begin{equation}\label{eq:average_power_cos}
	\tcbhighmath{
		\overline{P} = \frac{1}{2}\emf_m I_m \operatorname{Re}\left[e^{-i\phi_0}\right]
		= \frac{1}{2}\emf_m I_m \cos\phi_0.
	}
\end{equation}
Введемо ефективні значення амплітуд напруги і струму:
\begin{equation}\label{eq:effective_values}
	I_{\text{eff}} = \sqrt{\overline{I^2(t)}} = \frac{1}{\sqrt{2}}I_m, \qquad
	\emf_{\text{eff}} = \sqrt{\overline{\emf^2(t)}} = \frac{1}{\sqrt{2}}\emf_m.
\end{equation}
Тоді співвідношення \eqref{eq:average_power_cos} набирає вигляду
\begin{equation}\label{eq:average_power_effective}
	\tcbhighmath{
		\overline{P} = \emf_{\text{eff}} I_{\text{eff}} \cos\phi_0.
	}
\end{equation}
Ця формула показує, що чим більше зсув фаз між струмом і напругою, тим менше розсіювана в опорі потужність. У
випадку резонансу, коли $\tg\phi_0 = 0$, споживана потужність максимальна: $\overline{P} = \emf_{\text{eff}} I_{\text{eff}}$.

Перепишемо співвідношення \eqref{eq:average_power_complex}, враховуючи закон Ома: $\emf = IZ$, а
також рівності $\operatorname{Re} Z = R$, $|\hat{I}|^2 = I_m^2$:
\begin{equation}\label{eq:average_power_resistance}
	\tcbhighmath{
		\overline{P} = \frac{1}{2}\operatorname{Re}\left(\hat{\emf}(t)\hat{I}^*(t)\right)
		= \frac{1}{2}\operatorname{Re}\left(Z|\hat{I}|^2\right)
		= \frac{1}{2}RI_m^2 = RI_{\text{eff}}^2.
	}
\end{equation}
Остання формула показує, що за відсутності активного опору середня потужність змінного струму дорівнює нулю.

\begin{Historical}
	Формула $\overline{P} = RI_{\text{eff}}^2$ --- це, по суті, формула, яка розв'язала долю електричних мереж по всьому світу. Щоб доставити споживачеві задану потужність $P = U_{\text{eff}}I_{\text{eff}}\cos\phi_0$, можна взяти малу напругу й великий струм, або велику напругу й малий струм --- добуток лишається тим самим. Але втрати в дротах самої лінії залежать не від потужності, а саме від струму: $\overline{P}_{\text{втрат}} = RI_{\text{eff}}^2$. Тому, підвищивши напругу вдесятеро при тій самій потужності, ми зменшуємо струм у стільки ж разів, а втрати --- вже в сто разів. Саме на цьому \enquote{спіткнулася} постійнострумова мережа Едісона в 1880-х: підняти напругу постійного струму на відстані було технічно неможливо, тож електростанції доводилося ставити майже через кожні кілька кварталів. Джордж Вестінгауз і Нікола Тесла натомість запропонували систему змінного струму, в якій трансформатор дозволяє підняти напругу практично без втрат. У 1893 році саме цю систему обрали для гідроелектростанції на Ніагарському водоспаді (проєкти станції перед тим розглядала міжнародна комісія на чолі з лордом Кельвіном, який спершу схилявся до постійного струму); станція запрацювала 1895 року, а 1896 року вже передавала струм за кілька десятків кілометрів, до Буффало --- цю подію й вважають фактичним завершенням \enquote{війни струмів} на користь змінного струму. Відтоді магістральні лінії електропередачі (ЛЕП) тримають напругу в сотні кіловольт --- не з надмірності, а через оцю саму формулу $RI_{\text{eff}}^2$.
\end{Historical}

%% --------------------------------------------------------
\subsection*{Резонанс напруг і струмів}\label{sec:ac_resonance}
%% --------------------------------------------------------

У попередньому розділі ми вже зустрілися з явищем резонансу --- там воно постало як різке зростання амплітуди вимушених коливань при наближенні частоти ЕРС до власної частоти контуру. Тепер ми подивимося на це саме явище з боку імпедансів: виявиться, що резонанс --- це момент, коли реактивний опір кола обертається в нуль. Такий погляд дозволяє однаково просто описати і послідовний контур (резонанс напруг), і паралельний (резонанс струмів), а також поширити аналіз на довільні розгалужені кола.


\emphx{Резонансом напруг}\index{Резонанс!напруг} називають резонанс у послідовному
$LCR$-кон\-ту\-рі (рис.~\ref{tikz:LCR_circuit}) при зміні частоти коливань ЕРС. Згідно з \eqref{eq:Ohm_complex_derived} максимум амплітуди струму $I_m(\omega)$ досягається на частоті $\omega = \omega_0 = 1/\sqrt{LC}$ і становить $I_{m,\max} = \emf_m/R$ (рис.~\ref{tikz:current_resonance}). Це випливає з того, що на частоті $\omega_0$ індуктивний і ємнісний імпеданси точно компенсують один одного:
\begin{equation}\label{eq:impedance_compensation_2}
	Z_L + Z_C = \left(i\omega L + \frac{1}{i\omega C}\right)_{\omega = 1/\sqrt{LC}} = 0.
\end{equation}

\begin{SCfigure}[1][h!]
	\centering
	\localinput{current_resonance.tikz}
	\caption{Залежність амплітуди струму від частоти ЕРС у~послідовному $LCR$-контурі (рис.~\ref{tikz:LCR_circuit}): резонанс напруг.}
	\label{tikz:current_resonance}
\end{SCfigure}

В результаті імпеданс контуру зводиться до активного опору $R$, тож амплітуда струму досягає максимуму $I_{m,\max} = \emf_m/R$. Напруга на конденсаторі при цьому дорівнює:
\begin{equation*}
	V_{C,\max} = \frac{I_{m,\max}}{\omega_0 C} = \frac{\emf_m}{\omega_0 C R}.
\end{equation*}
Враховуючи, що $\omega_0 = 1/\sqrt{LC}$, маємо $\omega_0 C = \sqrt{C/L}$, звідки
\begin{equation*}
	V_{C,\max} = \frac{\emf_m}{R}\sqrt{\frac{L}{C}} = Q\,\emf_m,
\end{equation*}
де $Q = \frac{1}{R}\sqrt{\frac{L}{C}}$ --- добротність контуру \eqref{eq:Q_LCR}. Отже, напруга на конденсаторі в $Q$ разів перевищує напругу джерела. Те саме стосується й напруги на котушці. Саме тому це явище називають \emphx{резонансом напруг}.

\begin{Attention}
	Резонанс напруг --- це не лише корисне явище (виділення сигналу в радіоприймачах), а й потенційно \emph{небезпечне}. На резонансі напруги на котушці й конденсаторі хоч і рівні за модулем і протилежні за фазою (тому їхня сума дорівнює напрузі джерела), але \emph{кожна з них окремо} в $Q$ разів перевищує напругу джерела. Уже при $Q \sim 100$ і напрузі джерела $220$ В напруга на конденсаторі сягає $22\ 000$ В --- цього достатньо, щоб пробити ізоляцію конденсатора або спричинити розряд між витками котушки. У силовій електротехніці це явище називають \emphx{перенапругою} і враховують при проєктуванні обладнання; в радіотехніці ж його, навпаки, використовують для підсилення слабких сигналів.
\end{Attention}


\emphx{Резонансом струмів}\index{Резонанс!струмів} називають резонанс у паралельному $LCR$-контурі (рис.~\ref{tikz:parallel_circuit}), коли елементи $L$, $C$, $R$ з'єднані паралельно і підключені до джерела змінного струму. На відміну від резонансу напруг, тут амплітуда струму джерела вважається незмінною, тобто не залежить від параметрів елементів контуру.

\begin{SCfigure}[1.5][h!]
	\centering
	\localinput{parallel_circuit.tikz}
	\caption{Паралельний $LCR$-контур, що включає джерело заданого змінного струму $I(t) = I_m \cos\omega t$.}
	\label{tikz:parallel_circuit}
\end{SCfigure}

Оскільки елементи контуру з'єднані паралельно, імпеданс кола $Z$ визначається з рівності
\begin{equation}\label{eq:parallel_impedance}
	\frac{1}{Z} = \frac{1}{Z_R} + \frac{1}{Z_L} + \frac{1}{Z_C}, \qquad
	Z_R = R, \qquad Z_L = i\omega L, \qquad Z_C = \frac{1}{i\omega C}.
\end{equation}
Струм у зовнішній частині кола дорівнює сумі струмів через окремі елементи:
\begin{equation}\label{eq:current_sum}
	I = I_R + I_L + I_C.
\end{equation}
Позначимо напругу, створювану джерелом струму, як $U$. Тоді
\begin{equation}\label{eq:voltage_current}
	U = IZ.
\end{equation}
Оскільки напруга на всіх паралельно з'єднаних елементах однакова, то
\begin{equation}\label{eq:branch_currents}
	I_R = \frac{U}{Z_R} = \frac{U}{R}, \qquad
	I_L = \frac{U}{Z_L} = \frac{U}{i\omega L}, \qquad
	I_C = \frac{U}{Z_C} = i\omega C U.
\end{equation}
З урахуванням \eqref{eq:voltage_current} перепишемо ці вирази у вигляді
\begin{equation}\label{eq:branch_currents_Z}
	I_R = I\frac{Z}{R}, \qquad
	I_L = I\frac{Z}{Z_L}, \qquad
	I_C = I\frac{Z}{Z_C}.
\end{equation}

Амплітуда струму через активний опір дорівнює
\begin{equation}\label{eq:IR_amplitude}
	I_{R_0} = \frac{I_m}{\sqrt{1 + \left(\dfrac{R}{\omega L} - \omega C R\right)^2}}
	= \frac{I_m}{\sqrt{1 + \dfrac{4\gamma^2}{\omega^2}\left(1 - \dfrac{\omega^2}{\omega_0^2}\right)^2}}.
\end{equation}
Звідси видно, що $I_{R_0}$ має максимум при $\omega = \omega_0$ і становить $(I_{R0})_{\max} = I_m$. Це означає, що на резонансній частоті весь струм від джерела тече тільки через активний опір. Справа в тому, що опір блока з паралельно з'єднаних конденсатора й котушки на цій частоті обертається в нескінченність, і струм через нього не тече.

Амплітуди струмів через котушку й конденсатор на резонансній частоті дорівнюють
\begin{equation}\label{eq:resonance_currents_Q}
	I_{L0} = I_{C0} = QI_m,
\end{equation}
де
\begin{equation}\label{eq:Q_parallel}
	Q = \frac{R}{\omega_0 L} = \omega_0 R C = R\sqrt{\frac{C}{L}}
\end{equation}
--- добротність \emph{паралельного} контуру. Зверніть увагу, що вона відрізняється від добротності послідовного контуру \eqref{eq:Q_LCR}: там $Q = \frac{1}{R}\sqrt{L/C}$, тут --- $Q = R\sqrt{C/L}$. Фізично це означає, що в паралельному контурі добротність зростає зі збільшенням опору $R$ (менші втрати --- вища добротність), тоді як у послідовному --- навпаки, спадає.

Таким чином, на резонансній частоті по замкнутому $LC$-контуру циркулює струм, у~$Q$ разів більший за струм джерела, тобто тим більший, чим вища добротність кола.

Якщо замість джерела струму використовувати джерело ЕРС, то в зовнішньому колі на частоті $\omega_0$ спостерігався б \emph{мінімум} струму, а не максимум, як у послідовному контурі. Таке явище називають \emphx{антирезонансом}.

\begin{Attention}
	Резонанс струмів --- це дзеркальне відображення резонансу напруг. Якщо там різко зростали напруги на реактивних елементах, то тут різко зростають \emph{струми} в котушці й конденсаторі: на резонансі $I_{L0} = I_{C0} = QI_m$, тобто при високій добротності струм у $LC$-контурі в багато разів перевищує струм джерела. Це явище широко використовується в індукційних печах: паралельно котушці, що створює змінне магнітне поле, підключають ємність, і на резонансі струм у котушці зростає в багато разів, тоді як струм у підвідних дротах залишається відносно невеликим. Той самий принцип --- лише в мініатюрі --- лежить в основі побутових індукційних варильних поверхонь і безконтактної зарядки телефонів: замість тонни розплавленого металу тут вистачає котушки під склокерамічною плитою чи в зарядній панелі.

	У радіотехніці резонанс струмів використовують у приймальних пристроях для виділення слабкого сигналу у вузькому спектральному інтервалі поблизу резонансної частоти: паралельний $LC$-контур має на цій частоті великий опір, тому напруга на ньому максимальна, а сигнали інших частот пригнічуються.
\end{Attention}


%% --------------------------------------------------------
\subsection*{Правила Кірхгофа для змінного струму}
%% --------------------------------------------------------

Розрахунок\index{Кірхгофа правила!для змінного струму} довільних квазістаціонарних кіл (як нерозгалужених, так і розгалужених) при наявності змінних ЕРС може здійснюватися за правилами Кірхгофа, аналогічним тим, що мають місце у
випадку постійних струмів і напруг. Різниця полягає лише в заміні
опорів $R$ на імпеданси $Z$.

Правила такі:

\begin{enumerate}

	\item  Сума струмів, що входять у вузол (з урахуванням знаків), рівна нулю
	      (рис.~\ref{tikz:Kirchhoff}):
	      \begin{equation}\label{eq:Kirchhoff_current}
		      \tcbhighmath{
			      \sum_k I_k = 0.
		      }
	      \end{equation}
	      Дане правило є наслідком закону збереження заряду і стверджує, що у вузлі заряди не можуть накопичуватися.

	\item Для будь-якого замкнутого контуру в квазістаціонарному електричному колі виконується рівність
	      \begin{equation}\label{eq:Kirchhoff_voltage}
		      \sum_k I_k Z_k = \sum_i \emf_i,
	      \end{equation}
	      де $\emf_i$ --- ЕРС, що входять у вибраний контур, $I_k$ --- струми через елементи контуру, що мають імпеданси $Z_k$.

	      \begin{SCfigure}[1.75][h!]
		      \centering
		      \localinput{Kirchhoff.tikz}
		      \caption{Перше правило Кірхгофа: сума струмів, що входять у вузол, рівна нулю.}
		      \label{tikz:Kirchhoff}
	      \end{SCfigure}

	      Це правило безпосередньо випливає із закону індукції (теореми про циркуляцію для електричного поля в інтегральній формі):
	      \begin{equation}\label{eq:circulation_EMF}
		      \oint\limits_{\Gamma} E_l \, dl = -\frac{d\Phi}{dt}, \qquad \Phi = LI.
	      \end{equation}

\end{enumerate}

%%% --------------------------------------------------------
%\subsection*{Приклад розрахунку розгалуженого кола}
%%% --------------------------------------------------------

\begin{Example}{Приклад розрахунку розгалуженого кола}
Застосуємо правила Кірхгофа до конкретного розгалуженого кола (рис.~\ref{tikz:example_branched}): джерело ЕРС $\emf = \emf_m e^{i\omega t}$ навантажене на паралельно з'єднані резистор $R$ і конденсатор $C$, послідовно з якими ввімкнено котушку $L$. Позначимо вузол, у якому сходяться три вітки, буквою $A$, а другий вузол --- буквою $B$. Контур I містить джерело, котушку й резистор; контур II --- джерело, котушку й конденсатор.

\begin{center}%[h!]
	\centering
	\localinput{example_branched.tikz}
	\captionof{figure}{Розгалужене коло: котушка $L$ послідовно з паралельно з'єднаними $R$ і $C$. Вузли позначено $A$ і $B$; контури I та II показано умовно.}
	\label{tikz:example_branched}
\end{center}

Застосуємо \emphx{перше правило Кірхгофа} до вузла $A$: сума струмів, що виходять із вузла, дорівнює нулю. Струм $I$ приходить до вузла $A$, а струми $I_R$ і $I_C$ виходять із нього, тому
\begin{equation}\label{eq:example_KCL}
	I - I_R - I_C = 0
	\quad\Rightarrow\quad
	I = I_R + I_C.
\end{equation}

Застосуємо \emphx{друге правило Кірхгофа} до контуру I (джерело --- котушка --- резистор):
\begin{equation}\label{eq:example_KVL_I}
	Z_L I + R I_R = \emf.
\end{equation}
Для контуру II (джерело --- котушка --- конденсатор):
\begin{equation}\label{eq:example_KVL_II}
	Z_L I + Z_C I_C = \emf.
\end{equation}

Маємо систему трьох рівнянь \eqref{eq:example_KCL}--\eqref{eq:example_KVL_II} із трьома невідомими $I$, $I_R$, $I_C$. Розв'яжемо її. З \eqref{eq:example_KVL_I} і \eqref{eq:example_KVL_II} виразимо $I_R$ і $I_C$ через $I$:
\begin{equation*}
	I_R = \frac{\emf - Z_L I}{R}, \qquad
	I_C = \frac{\emf - Z_L I}{Z_C}.
\end{equation*}
Підставимо ці вирази в \eqref{eq:example_KCL}:
\begin{equation*}
	I = \frac{\emf - Z_L I}{R} + \frac{\emf - Z_L I}{Z_C}.
\end{equation*}
Розв'язуючи відносно $I$, отримуємо
\begin{equation}\label{eq:example_current}
	I = \frac{\emf}{Z_L + \dfrac{1}{\dfrac{1}{R} + \dfrac{1}{Z_C}}}.
\end{equation}
Вираз у знаменнику --- це повний імпеданс кола $Z$. Дійсно, паралельна ділянка $R \parallel C$ має імпеданс
\begin{equation}\label{eq:example_parallel}
	Z_{RC} = \frac{1}{\dfrac{1}{R} + \dfrac{1}{Z_C}} = \frac{R}{1 + i\omega RC},
\end{equation}
а повний імпеданс кола, згідно з правилом послідовного складання, дорівнює
\begin{equation}\label{eq:example_total}
	Z = Z_L + Z_{RC} = i\omega L + \frac{R}{1 + i\omega RC}.
\end{equation}
Підставляючи \eqref{eq:example_parallel} і \eqref{eq:example_total} у \eqref{eq:example_current}, отримуємо остаточно
\begin{equation}\label{eq:example_current_explicit}
	I = \frac{\emf_m e^{i\omega t}}{i\omega L + \dfrac{R}{1 + i\omega RC}}
	  = \emf_m e^{i\omega t}\,\frac{1 + i\omega RC}{R(1 - \omega^2 LC) + i\omega L}.
\end{equation}
Напруга на паралельній ділянці:
\begin{equation*}
	U_{RC} = I\,Z_{RC}
	       = \frac{\emf_m e^{i\omega t}\,R}{R(1 - \omega^2 LC) + i\omega L}.
\end{equation*}
Струми через резистор і конденсатор:
\begin{equation*}
	I_R = \frac{U_{RC}}{R} = \frac{\emf_m e^{i\omega t}}{R(1 - \omega^2 LC) + i\omega L},
\end{equation*}
\begin{equation*}
	I_C = i\omega C\,U_{RC} = \frac{i\omega RC\,\emf_m e^{i\omega t}}{R(1 - \omega^2 LC) + i\omega L}.
\end{equation*}

Перевіримо виконання першого правила Кірхгофа \eqref{eq:example_KCL}:
\begin{equation*}
	I_R + I_C = \frac{\emf_m e^{i\omega t}(1 + i\omega RC)}{R(1 - \omega^2 LC) + i\omega L} = I,
\end{equation*}
що збігається з \eqref{eq:example_current_explicit}. Отже, розв'язок задовольняє обидва правила Кірхгофа.

Із отриманих виразів видно, як розгалужене коло розраховується чисто алгебраїчно: спочатку записуємо рівняння Кірхгофа через імпеданси, потім розв'язуємо систему відносно невідомих струмів. Саме в цьому й полягає практична сила методу комплексних амплітуд: диференціальні рівняння замінюються системою лінійних алгебраїчних рівнянь із комплексними коефіцієнтами.
\end{Example}

%=========================================================
\section{Трансформатор}\hypertarget{transformer}{}
\label{sec:transformer}
%=========================================================

Принцип дії \emphx{трансформатора}\index{Трансформатор}~--- взаємну індукцію двох обмоток на спільному осерді~--- ми вже
обговорили в~розділі~\ref{MagEnergy} (рис.~\ref{pic:Transformer}), де коефіцієнт трансформації отримано із закону
збереження енергії для ідеального зв'язку обмоток. Тут розглянемо трансформатор як систему двох зв'язаних контурів змінного
струму (рис.~\ref{pic:transformer_circuit}): обмотки мають опори $R_1$, $R_2$, власні індуктивності $L_1$, $L_2$
і~взаємну індуктивність $L_{12}$ (у~СІ $\Phi_1=L_1I_1+L_{12}I_2$, $\Phi_2=L_{12}I_1+L_2I_2$), а~вторинну обмотку навантажено
імпедансом $Z_{\text{н}}$. Коефіцієнт трансформації тут зручно означити як $n=N_1/N_2$ (у~розділі~\ref{MagEnergy}
$\eta=N_2/N_1=1/n$).

\begin{figure}[h!]
	\centering
	\localinput{transformer_circuit.tikz}
	\caption{Еквівалентна схема трансформатора: котушки $L_1$, $L_2$ з~крапками фазування (вказують на однакову полярність
		виводів), опори $R_1$, $R_2$, взаємна індуктивність $L_{12}$ та навантаження~$Z_{\text{н}}$.}
	\label{pic:transformer_circuit}
\end{figure}

Друге правило Кірхгофа для комплексних амплітуд дає
\begin{align}
	\hat{\emf}_1 &= \left(R_1 + i\omega L_1\right)\hat{I}_1 + i\omega L_{12} \hat{I}_2,
	\label{eq:transformer_complex_1} \\[4pt]
	0 &= i\omega L_{12} \hat{I}_1 + \left(R_2 + i\omega L_2 + Z_{\text{н}}\right)\hat{I}_2
	\label{eq:transformer_complex_2}
\end{align}
(у~вторинному колі стороннього джерела немає: струм $\hat{I}_2$ збуджує лише ЕРС взаємної індукції). Виключивши $\hat{I}_2$,
знаходимо \emphx{вхідний імпеданс} трансформатора
\begin{equation}\label{eq:input_impedance}
	\tcbhighmath{
		Z_{\text{вх}} = \frac{\hat{\emf}_1}{\hat{I}_1} = R_1 + i\omega L_1 + \frac{(\omega L_{12})^2}{R_2 + i\omega L_2 + Z_{\text{н}}};
	}
\end{equation}
останній доданок називають \emphx{внесеним опором}\index{Трансформатор!внесений опір}: через нього джерело первинного кола
\enquote{відчуває} навантаження вторинного.

\emph{Ідеальний трансформатор.}\index{Трансформатор!ідеальний} Нехай опори обмоток малі ($R_1=R_2=0$), потік не
розсіюється ($L_{12}=\sqrt{L_1L_2}$, $L_1/L_2=(N_1/N_2)^2=n^2$), а~індуктивності великі ($\omega L_2\gg|Z_{\text{н}}|$~---
для цього й~потрібне феромагнітне осердя). Тоді з~\eqref{eq:transformer_complex_1}, \eqref{eq:transformer_complex_2} для
напруг на затискачах обмоток $\hat U_1=\hat{\emf}_1$, $\hat U_2=Z_{\text{н}}\hat I_2$ отримуємо
\begin{equation}\label{eq:ideal_transformer}
	\tcbhighmath{
		\frac{|\hat{U}_1|}{|\hat{U}_2|} = n, \qquad \frac{|\hat{I}_1|}{|\hat{I}_2|} = \frac1n, \qquad Z_{\text{вх}} = n^2 Z_{\text{н}}.
	}
\end{equation}
У~скільки разів трансформатор підвищує напругу, у~стільки ж разів знижує струм, а~потужність
$|\hat U_1||\hat I_1|=|\hat U_2||\hat I_2|$ зберігається. Остання формула показує, що трансформатор \enquote{перераховує}
опір навантаження в~первинне коло в~$n^2$ разів; цим користуються для \emphx{узгодження навантаження} з~джерелом
(задача~\ref{prb:MatchingSpeaker}).

Головне застосування трансформатора~--- передача електроенергії. Втрати в~лінії з~опором $R_\text{л}$ при переданій
потужності $P$ дорівнюють $R_\text{л}I^2 = R_\text{л}\left(P/U\cos\phi\right)^2$, тобто обернено пропорційні квадрату
напруги: підвищивши напругу в~$n$ разів, зменшуємо втрати в~$n^2$ разів (задача~\ref{prb:TransmissionLoss}). Тому магістральні
лінії працюють на сотнях кіловольт, а~біля споживачів напругу знову знижують.

%% --------------------------------------------------------
\section{Частотні фільтри}\hypertarget{frequency-filters}{}
%% --------------------------------------------------------

Оскільки реактивні опори залежать від частоти ($X_L = \omega L$, $X_C = 1/(\omega C)$), кола з реактивними елементами по-різному пропускають сигнали різних частот. Ця властивість лежить в основі \emphx{частотних фільтрів}\index{Фільтри частотні} --- пристроїв, що виділяють або пригнічують певні ділянки спектра.

%% --------------------------------------------------------
\subsection*{Фільтр низьких частот (ФНЧ)}
%% --------------------------------------------------------

Найпростіший ФНЧ --- це $RC$-коло, в якому вихідна напруга знімається з конденсатора (рис.~\ref{tikz:lowpass}). Коефіцієнт передачі за напругою:
\begin{equation}\label{eq:lowpass_transfer}
	K(\omega) = \frac{U_{\text{вих}}}{U_{\text{вх}}}
	          = \frac{Z_C}{R + Z_C}
	          = \frac{1}{1 + i\omega RC}.
\end{equation}

\begin{SCfigure}[1][h!]
	\centering
	\localinput{lowpass.tikz}
	\caption{Фільтр низьких частот: вихідна напруга знімається з конденсатора $C$.}
	\label{tikz:lowpass}
\end{SCfigure}

Модуль коефіцієнта передачі:
\begin{equation}\label{eq:lowpass_magnitude}
	|K(\omega)| = \frac{1}{\sqrt{1 + (\omega RC)^2}}.
\end{equation}
На низьких частотах $\omega \ll 1/(RC)$ маємо $|K| \approx 1$ --- сигнал проходить без ослаблення. На високих частотах $\omega \gg 1/(RC)$ маємо $|K| \approx 1/(\omega RC) \to 0$ --- сигнал пригнічується. Гранична частота, на якій $|K| = 1/\sqrt{2}$, дорівнює
\begin{equation}\label{eq:lowpass_cutoff}
	\omega_c = \frac{1}{RC}.
\end{equation}
Такий фільтр використовується, наприклад, для згладжування пульсацій у випрямлячах живлення: конденсатор великої ємності шунтує високочастотні складові, залишаючи майже сталу напругу.

%% --------------------------------------------------------
\subsection*{Фільтр високих частот (ФВЧ).}
%% --------------------------------------------------------

Якщо в тому самому $RC$-колі знімати вихідну напругу з резистора (рис.~\ref{tikz:highpass}), отримаємо ФВЧ:
\begin{equation}\label{eq:highpass_transfer}
	K(\omega) = \frac{R}{R + Z_C} = \frac{i\omega RC}{1 + i\omega RC},
	\qquad
	|K(\omega)| = \frac{\omega RC}{\sqrt{1 + (\omega RC)^2}}.
\end{equation}
Тепер на низьких частотах $|K| \to 0$, а на високих $|K| \to 1$. Такий фільтр застосовується, наприклад, для розділення частот у гучномовцях: високочастотний динамік (твітер) підключається через конденсатор, який не пропускає низькочастотні складові.

\begin{SCfigure}[1][h!]
	\centering
	\localinput{highpass.tikz}
	\caption{Фільтр високих частот: вихідна напруга знімається з резистора $R$.}
	\label{tikz:highpass}
\end{SCfigure}

%% --------------------------------------------------------
\subsection*{Смуговий фільтр.}
%% --------------------------------------------------------
Послідовний $LCR$-контур, у якому вихідна напруга знімається з резистора (рис.~\ref{tikz:bandpass}), є смуговим фільтром. Його коефіцієнт передачі
\begin{equation}\label{eq:bandpass_transfer}
	K(\omega) = \frac{R}{R + i(\omega L - 1/\omega C)}
	          = \frac{1}{1 + iQ\left(\dfrac{\omega}{\omega_0} - \dfrac{\omega_0}{\omega}\right)},
\end{equation}
де $Q = \omega_0 L/R$ --- добротність. Максимум $|K| = 1$ досягається на резонансній частоті $\omega_0 = 1/\sqrt{LC}$, а ширина смуги пропускання $\Delta\omega = \omega_0/Q$. Чим вища добротність, тим вужча смуга --- саме тому високодобротні контури використовують для виділення одного радіоканалу з безлічі інших.

\begin{figure}[h!]
	\centering
	\localinput{bandpass.tikz}
	\caption{Смуговий фільтр на основі послідовного $LCR$-контуру: вихідна напруга знімається з резистора $R$.}
	\label{tikz:bandpass}
\end{figure}

Порівняти поведінку всіх трьох фільтрів зручно на одному графіку амплітудно-частотних характеристик (рис.~\ref{tikz:filter_responses}).

\begin{figure}[h!]
	\centering
	\localinput{filter_responses.tikz}
	\caption{Амплітудно-частотні характеристики трьох фільтрів: ФНЧ (синій), ФВЧ (червоний) і смугового (зелений). Для смугового прийнято $Q = 5$.}
	\label{tikz:filter_responses}
\end{figure}

\begin{Attention}
	Частотні фільтри --- це, мабуть, найпоширеніший застосунок кіл змінного струму. У радіоприймачі вхідний $LC$-контур виділяє потрібну станцію з усього ефіру; у підсилювачі звуку $RC$-ланцюжки розділяють смуги між високочастотним і низькочастотним динаміками; у блоці живлення комп'ютера конденсатори й дроселі придушують високочастотні завади, щоб вони не проникали в мережу; у цифрових схемах фільтри прибирають шум із сигналів перед оцифровуванням. Спільна ідея в усіх випадках одна: реактивний опір залежить від частоти, і цим можна керувати тим, які частоти \enquote{пройдуть}, а які --- ні.
\end{Attention}

%% --------------------------------------------------------
\section{Трифазний струм}\hypertarget{three-phase-current}{}
%% --------------------------------------------------------

Для передачі великих потужностей і~живлення двигунів використовують \emphx{трифазні системи}\index{Трифазний струм}.
Генератор має три однакові обмотки статора, зсунуті одна відносно одної на $120^\circ$; магніт ротора, що обертається
з~кутовою швидкістю $\omega$, наводить у~них ЕРС однакової амплітуди, зсунуті в~часі на третину періоду
(рис.~\ref{pic:ThreePhaseGenerator}):
\begin{align}
	\emf_1(t) &= \emf_m \sin(\omega t), \notag \\
	\emf_2(t) &= \emf_m \sin\left(\omega t - \frac{2\pi}{3}\right), \label{eq:emf_system} \\
	\emf_3(t) &= \emf_m \sin\left(\omega t + \frac{2\pi}{3}\right). \notag
\end{align}

\begin{figure}[h!]
	\centering
	\localinput{ThreePhaseGenerator.tikz}
	\caption{Схематичне зображення трифазного генератора}
	\label{pic:ThreePhaseGenerator}
\end{figure}

Сума трьох ЕРС у~будь-який момент дорівнює нулю. Тому, з'єднавши кінці обмоток в~одну точку (\enquote{зірка}), енергію
можна передавати трьома дротами замість шести: за однакового навантаження фаз струм у~спільному (нейтральному) дроті
дорівнює нулю. Напруга між двома лінійними дротами в~$\sqrt3$ разів більша за напругу однієї фази (звідси пара $380$
і~$220$~В). Таку систему 1889~р. запропонував М.~Доливо-Добровольський, а~1891~р. її переваги показала передача енергії
лінією Лауффен--Франкфурт завдовжки $175$~км.

\emph{Обертове магнітне поле.}\index{Обертове магнітне поле} Три однакові котушки, осі яких розташовані під кутом
$120^\circ$ одна до одної, живлять трифазними струмами $I_k = I_m\sin\left(\omega t - 2\pi(k-1)/3\right)$. Кожна створює
пульсуюче поле вздовж своєї осі, $\vect B_k = B_m\sin\left(\omega t - 2\pi(k-1)/3\right)\vect e_k$, а~їхня сума
\begin{equation}
	\vect B = \vect B_1 + \vect B_2 + \vect B_3, \qquad |\vect B| = \frac32 B_m = \const,
\end{equation}
обертається з~кутовою швидкістю $\omega$ (задача~\ref{prb:RotatingField}). Таке поле без жодних рухомих частин захоплює
провідний ротор (струми Фуко, розділ~\ref{ElMagInduction})~--- на цьому працюють асинхронні двигуни.

\emph{Стала потужність.} За однакового навантаження фаз сума миттєвих потужностей трьох фаз не залежить від часу:
$p(t) = 3U_\text{ф}I_\text{ф}\cos\phi = \const$, тоді як в~однофазному колі потужність пульсує з~частотою $2\omega$. Це
забезпечує рівномірну роботу потужних генераторів і~двигунів.


%% --------------------------------------------------------
\section*{Підсумки}
\phantomsection
\addcontentsline{toc}{section}{Підсумки}
%% --------------------------------------------------------

Електричне коло можна описувати простими співвідношеннями (закон Ома, правила Кірхгофа) без повних рівнянь Максвелла, якщо виконується умова \emphx{квазістаціонарності}: довжина кола значно менша за довжину хвилі, $l \ll \lambda = cT$ (або $\tau \sim l/c \ll T$). Тоді коло складається із зосереджених елементів $R$, $L$, $C$, і процеси в ньому описує рівняння $L\dot I + IR + q/C = \emf$. За відсутності ЕРС воно перетворюється на рівняння вільних коливань $\ddot q + 2\gamma\dot q + \omega_0^2 q = 0$, де $\gamma = R/2L$, $\omega_0 = 1/\sqrt{LC}$. При $R = 0$ коливання гармонічні з періодом $T = 2\pi\sqrt{LC}$ (формула Томсона), причому енергія періодично перетікає між електричним полем конденсатора й магнітним полем котушки. За слабкого згасання ($\gamma < \omega_0$) амплітуда спадає як $e^{-\gamma t}$ з частотою $\omega = \sqrt{\omega_0^2 - \gamma^2}$, а при $\gamma \geqslant \omega_0$ коливання не виникають. Швидкість згасання характеризують логарифмічний декремент $\delta = \gamma T$ та добротність $Q = \omega/2\gamma = \pi/\delta = 2\pi W/\Delta W$; для послідовного $LCR$-контуру $Q = \frac1R\sqrt{L/C}$.

Під дією гармонічної ЕРС, після загасання вільних коливань, у контурі лишаються \emphx{вимушені} коливання на частоті джерела. Їх зручно шукати методом комплексних амплітуд, який замінює диференціальне рівняння алгебраїчним. Амплітуда напруги на конденсаторі має максимум (\emphx{резонанс}) при $\omega_m = \sqrt{\omega_0^2 - 2\gamma^2}$; за слабкого згасання $V_{\max} = QV_{\text{ст}}$, а ширина резонансної кривої по рівню $1/\sqrt2$ дорівнює $\Delta\omega = 2\gamma = \omega_0/Q$: чим вища добротність, тим вужчий і вищий резонанс. Те саме явище з боку імпедансів виглядає так: закон Ома набуває форми $\hat I_m Z = \emf_m$, де \emphx{імпеданс} $Z = R + i(\omega L - 1/\omega C)$, а окремі елементи мають $Z_R = R$, $Z_L = i\omega L$, $Z_C = 1/(i\omega C)$. Правила Кірхгофа зберігають форму, якщо замінити $R$ на $Z$. Середня потужність $\overline{P} = \emf_\text{еф}I_\text{еф}\cos\phi_0 = RI_\text{еф}^2$ визначається лише активним опором, тому електроенергію вигідно передавати при високій напрузі. У послідовному контурі при резонансі напруг напруги на $L$ і $C$ у $Q$ разів перевищують напругу джерела; у паралельному контурі при резонансі струмів струм у гілках $L$ і $C$ у~$Q$ разів перевищує струм джерела, а $Q = R\sqrt{C/L}$.

Метод імпедансів дозволяє розрахувати й практично важливі пристрої. \emphx{Трансформатор} як систему двох магнітно зв'язаних контурів описують рівняння з взаємною індуктивністю $L_{12}$; для ідеального трансформатора $|U_1/U_2| = N_1/N_2 = n$, $|I_1/I_2| = 1/n$ і $Z_\text{вх} = n^2Z_\text{н}$, що дає змогу узгоджувати навантаження з джерелом. \emphx{Частотні фільтри} використовують залежність реактивних опорів від частоти: для $RC$-ланцюжка гранична частота $\omega_c = 1/RC$, а смуговий фільтр має смугу пропускання $\Delta\omega = \omega_0/Q$. У \emphx{трифазній системі} три ЕРС зсунуті на $120^\circ$ і в сумі дають нуль, тож енергію передають трьома дротами; трифазний струм створює обертове магнітне поле сталої величини й забезпечує сталу миттєву потужність.

Уся ця теорія трималася на умові $l \ll \lambda$. Коли довжина лінії стає порівнянною з довжиною хвилі, струм і напруга вже не однакові вздовж кола, і потрібна теорія \emph{довгих ліній}, з якої ми почнемо наступний розділ (розділ~\ref{Telegraph}).


%% --------------------------------------------------------
\section*{Задачі}
\addcontentsline{toc}{section}{Задачі}
%% --------------------------------------------------------

% задачі цього розділу короткі: не розриваємо їх між сторінками
\begingroup
\tcbset{answersolutionstyle/.append style={unbreakable}}

\begin{problem}\label{prb:QuasiStationarity}
	\emph{Квазістаціонарність.} При якій найбільшій частоті коло довжиною $3$~м ще можна вважати квазістаціонарним, якщо вимагати $l\leqslant\lambda/10$? Яка частка довжини хвилі мережі 50~Гц припадає на ЛЕП завдовжки $1000$~км?
	\begin{solution}
		$\nu_{\max}=c/(10l)=10$~МГц; $l/\lambda=10^{6}\ \text{м}/(6\cdot10^{6}\ \text{м})\approx0{,}17$ --- лінія вже становить відчутну частку довжини хвилі, і квазістаціонарне наближення неприпустиме.
	\end{solution}
\end{problem}

\begin{problem}\label{prb:FreeOscillationsNumbers}
	\emph{Вільні коливання в контурі.} Контур містить $L=0{,}5$~мГн, $C=200$~пФ і $R=5$~Ом. Знайдіть власну частоту, період, коефіцієнт згасання, логарифмічний декремент, добротність, кількість коливань, за які амплітуда спадає в $e$ разів, та критичний опір, при якому коливання стають неможливими.
	\begin{solution}
		$\omega_0=1/\sqrt{LC}\approx3{,}16\cdot10^{6}$~с$^{-1}$ ($\nu\approx503$~кГц), $T\approx1{,}99$~мкс; $\gamma=R/2L=5\cdot10^{3}$~с$^{-1}$; $\delta=\gamma T\approx9{,}9\cdot10^{-3}$; $Q=\dfrac{1}{R}\sqrt{L/C}\approx316$; $n=1/\delta\approx100$; $R_\text{кр}=2\sqrt{L/C}\approx3{,}2$~кОм.
	\end{solution}
\end{problem}

\begin{problem}\label{prb:ForcedResonanceNumbers}
	\emph{Резонанс у послідовному контурі.} До контуру з попередньої задачі приєднано генератор із $\emf_m=1$~В. Знайдіть амплітуду струму й напруги на конденсаторі на резонансі та ширину резонансної кривої.
	\begin{solution}
		$I_{m,\max}=\emf_m/R=0{,}2$~А; $V_{C,\max}=Q\emf_m\approx316$~В (у $316$ разів більше за напругу джерела); $\Delta\omega=2\gamma=10^{4}$~с$^{-1}$, тобто $\Delta\nu=\omega_0/(2\pi Q)\approx1{,}6$~кГц.
	\end{solution}
\end{problem}

\begin{problem}\label{prb:ImpedanceNumbers}
	\emph{Імпеданс і потужність.} Послідовно з'єднані резистор $R=100$~Ом, котушка $L=0{,}1$~Гн та конденсатор $C=10$~мкФ підключені до мережі $220$~В, $50$~Гц. Знайдіть імпеданс, зсув фаз, ефективний струм і споживану потужність. Чим викликано малий $\cos\phi$?
	\begin{solution}
		$X_L=31{,}4$~Ом, $X_C=318$~Ом, $Z=100-287i$~Ом, $|Z|\approx304$~Ом; $\phi=\arg Z\approx-71^\circ$ (струм випереджає напругу, коло ємнісне); $I_\text{еф}=U/|Z|\approx0{,}72$~А; $P=I_\text{еф}^2R\approx52$~Вт; $\cos\phi=R/|Z|\approx0{,}33$ --- через великий реактивний опір $|X_L-X_C|$.
	\end{solution}
\end{problem}

\begin{problem}\label{prb:ParallelResonanceNumbers}
	\emph{Резонанс струмів.} Паралельний контур ($R=100$~кОм, $L=1$~мГн) настроєно на $1$~МГц і під'єднано до джерела струму з амплітудою $I_m=1$~мА. Знайдіть ємність, добротність, амплітуду напруги на контурі та амплітуди струмів у котушці й конденсаторі.
	\begin{solution}
		$C=1/(\omega_0^2L)\approx25$~пФ; $Q=R/(\omega_0L)\approx15{,}9$; $U_m=I_mR=100$~В; $I_{L0}=I_{C0}=QI_m\approx16$~мА --- у $Q$ разів більше за струм джерела.
	\end{solution}
\end{problem}

\begin{problem}\label{prb:PowerFactor}
	\emph{Коефіцієнт потужності.} Підприємство споживає активну потужність $P=1$~МВт при $\cos\phi=0{,}6$ (індуктивне навантаження) від однофазного джерела $10$~кВ, $50$~Гц. Знайдіть струм, реактивну потужність і ємність батареї конденсаторів, що доводить $\cos\phi$ до $1$. Як зміняться втрати в лінії?
	\begin{solution}
		$I=P/(U\cos\phi)\approx167$~А; $Q_\text{р}=P\tg\phi\approx1{,}33$~Мвар; батарея, що створює $Q_C=U^2\omega C=Q_\text{р}$, має $C=Q_\text{р}/(\omega U^2)\approx42$~мкФ. Після компенсації $I=P/U=100$~А, втрати $\propto I^2$ зменшуються в $(0{,}6)^{-2}\approx2{,}8$ раза.
	\end{solution}
\end{problem}

\begin{problem}\label{prb:MatchingSpeaker}
	\emph{Узгодження за допомогою трансформатора.} Підсилювач є джерелом із ЕРС $80$~В (ефективна) і внутрішнім опором $800$~Ом, гучномовець має опір $8$~Ом. Знайдіть потужність у гучномовці при безпосередньому підключенні та через ідеальний трансформатор, що узгоджує опори. Який коефіцієнт трансформації $n=N_1/N_2$ потрібен?
	\begin{solution}
		без трансформатора $P=\dfrac{\emf^2R_\text{н}}{(R_\text{дж}+R_\text{н})^2}\approx0{,}078$~Вт. Узгодження: $n=\sqrt{R_\text{дж}/R_\text{н}}=10$; напруга на первинній обмотці $\emf/2=40$~В, на гучномовці $U_2=U_1/n=4$~В, $P=U_2^2/R_\text{н}=2$~Вт --- у $\approx25$ разів більше.
	\end{solution}
\end{problem}

\begin{problem}\label{prb:TransmissionLoss}
	\emph{Втрати в лінії.} Потужність $100$~МВт передають лінією з сумарним опором проводів $20$~Ом при напрузі $110$~кВ і $500$~кВ ($\cos\phi=1$, розрахунок ведіть як для однофазної лінії). Порівняйте втрати.
	\begin{solution}
		$P_\text{втр}=R(P/U)^2$: $\approx16{,}5$~МВт ($16{,}5\,\%$) при $110$~кВ і $0{,}8$~МВт ($0{,}8\,\%$) при $500$~кВ --- у $(500/110)^2\approx21$ раз менше.
	\end{solution}
\end{problem}

\begin{problem}\label{prb:RCFilter}
	\emph{$RC$-фільтр низьких частот.} $R=1$~кОм, $C=0{,}16$~мкФ. Знайдіть граничну частоту, а також модуль і фазу коефіцієнта передачі при $10$~кГц. У скільки разів (і на скільки децибел) послаблюється сигнал?
	\begin{solution}
		$\nu_c=1/(2\pi RC)\approx1$~кГц; $\omega RC\approx10$, $|K|=1/\sqrt{1+(\omega RC)^2}\approx0{,}10$, $\arg K=-\arctan(\omega RC)\approx-84^\circ$; послаблення у $\approx10$ разів, тобто $\approx20$~дБ.
	\end{solution}
\end{problem}

\begin{problem}\label{prb:RotatingField}
	\emph{Обертове поле.} Покажіть, що при зміні місцями двох фаз у котушках трифазної системи напрям обертання магнітного поля змінюється на протилежний.
	\begin{solution}
		при $I_2\leftrightarrow I_3$ маємо $B_y=+\frac32B_m\cos\omega t$ при тому ж $B_x=\frac32B_m\sin\omega t$, тому кут $\theta$ вектора $\mathbf{B}$ визначається рівностями $\cos\theta=\sin\omega t$, $\sin\theta=\cos\omega t$, тобто $\theta=90^\circ-\omega t$ і вектор обертається за годинниковою стрілкою; раніше $\theta=\omega t-90^\circ$ (проти годинникової).
	\end{solution}
\end{problem}
\endgroup


\begin{Questions}
\begin{quizlist}
	\item Що таке квазістаціонарний процес? Сформулюйте умову квазістаціонарності через довжину кола та довжину хвилі. Чому для будинкової проводки при 50~Гц вона виконується з великим запасом, а для довгої ЛЕП це може бути не так?
	\item Чим система із зосередженими елементами відрізняється від системи з розподіленими параметрами (довгої лінії)? Для чого можна застосовувати закон Ома та правила Кірхгофа?
	\item Виведіть рівняння $L\dot I + IR + q/C = \emf$ для послідовного $LCR$-контуру (рис.~\ref{tikz:LCR_circuit}). Яку роль відіграє теорема про циркуляцію електричного поля і як вона враховує ЕРС індукції в котушці?
	\item Запишіть рівняння вільних коливань $\ddot q + 2\gamma\dot q + \omega_0^2 q = 0$. Що означають $\gamma$ та $\omega_0$ і як вони виражаються через $R$, $L$, $C$?
	\item Розв'яжіть рівняння вільних коливань при $R=0$. Запишіть формулу Томсона. Чим визначаються амплітуда та початкова фаза коливань?
	\item Як знаходять розв'язок рівняння згасаючих коливань за допомогою комплексних чисел? Які три випадки виникають залежно від співвідношення $\omega_0$ та $\gamma$ і який вигляд має розв'язок у кожному з них?
	\item Що таке коефіцієнт згасання та час згасання? Як виглядає графік згасаючих коливань (рис.~\ref{tikz:damped}) і який вигляд має закон зміни амплітуди?
	\item Виведіть рівняння енергетичного балансу $\frac{d}{dt}\left(\frac{LI^2}{2}+\frac{q^2}{2C}\right) = -RI^2$. Що означають його доданки?
	\item Як перетворюється енергія в контурі без втрат? Чому енергії $W_L$ та $W_C$ коливаються з однаковою амплітудою в протифазі, а їхня сума залишається сталою? Який зсув фаз між $q(t)$ і $I(t)$?
	\item Що таке логарифмічний декремент згасання? Виведіть його вираз через $\gamma$ та $T$ і поясніть зв'язок з числом коливань $n$, за які амплітуда спадає в $e$ разів.
	\item Що таке добротність коливальної системи? Запишіть її через $\omega$ та $\gamma$, через логарифмічний декремент та через відношення енергії до втрат за період. Чому дорівнює добротність послідовного $LCR$-контуру?
	\item Яких порядків величини досягає добротність у звичайному $LC$-контурі, кварцовому резонаторі та надпровідному резонаторі прискорювача? Що це означає для тривалості \enquote{дзвону} системи?
	\item У чому полягає метод комплексних амплітуд? Що таке комплексна амплітуда і як з її модуля та аргумента отримати амплітуду та початкову фазу коливання?
	\item Що таке частотна характеристика контуру? Виведіть вираз для амплітуди $V_m$ та фазового зсуву вимушених коливань напруги на конденсаторі.
	\item Опишіть амплітудно-частотну (рис.~\ref{tikz:amplitude_response}) та фазово-частотну (рис.~\ref{tikz:phase_response}) характеристики контуру. Як змінюється їхній вигляд зі збільшенням згасання?
	\item Що таке резонанс? Чому дорівнює резонансна частота $\omega_m$ при $\gamma\ne0$ і за якої умови резонанс відсутній?
	\item Виведіть значення $V_{\max}$ при слабкому згасанні та поясніть, чому відношення $V_{\max}/V_{\text{ст}}$ дорівнює добротності. Що таке статична амплітуда $V_{\text{ст}}$?
	\item Виведіть ширину резонансної кривої по рівню $1/\sqrt2$ і запишіть її через добротність. Як зв'язані добротність та вузькість резонансу?
	\item Що таке змінний струм? Чим закон Ома в комплексній формі відрізняється від закону Ома для постійного струму?
	\item Що таке імпеданс? Запишіть імпеданси $R$, $L$, $C$ та послідовного $LCR$-кола. Чим активний опір відрізняється від реактивного?
	\item Який зсув фаз між струмом і напругою на резисторі, конденсаторі та котушці? Чому множення на $i$ у комплексній площині означає поворот на $\pi/2$?
	\item Як побудувати векторну діаграму послідовного $LCR$-контуру (рис.~\ref{tikz:vector_diagram_full})? Що можна визначити з неї, не виконуючи обчислень?
	\item Виведіть формулу для середньої потужності змінного струму $\overline P = \frac12\emf_mI_m\cos\phi_0$. Що таке ефективні значення струму й напруги і чому потужність залежить від зсуву фаз?
	\item Чому середня потужність змінного струму дорівнює нулю за відсутності активного опору? Чому при тій самій потужності вигідно передавати електроенергію при високій напрузі?
	\item Що таке резонанс напруг? Чому в резонансі імпеданс контуру дорівнює $R$, а напруга на конденсаторі та котушці в $Q$ разів перевищує напругу джерела? Чим небезпечний цей ефект?
	\item Що таке резонанс струмів у паралельному контурі (рис.~\ref{tikz:parallel_circuit})? Як записується імпеданс паралельного кола і що відбувається зі струмами в гілках на резонансній частоті?
	\item Чим відрізняється добротність паралельного контуру від добротності послідовного? Що таке антирезонанс?
	\item Сформулюйте правила Кірхгофа для змінного струму. Чим вони відрізняються від правил для постійного струму і як пов'язані із законом збереження заряду та законом індукції?
	\item Як розраховується розгалужене коло методом комплексних амплітуд? Які кроки включає розв'язання прикладу з $L$, $R$, $C$ (рис.~\ref{tikz:example_branched})?
	\item Чому трансформатор працює лише на змінному струмі?
	\item Запишіть рівняння трансформатора як двох зв'язаних контурів \eqref{eq:transformer_complex_1}, \eqref{eq:transformer_complex_2}. Що таке внесений опір?
	\item Виведіть співвідношення ідеального трансформатора для напруг, струмів і~вхідного імпедансу. Навіщо узгоджують навантаження з~джерелом?
	\item Що таке частотний фільтр? Виведіть коефіцієнт передачі та граничну частоту $RC$-фільтра низьких частот (рис.~\ref{tikz:lowpass}). Як працює смуговий фільтр на основі $LCR$-контуру і~від чого залежить ширина смуги?
	\item Чому сума трьох ЕРС симетричної трифазної системи дорівнює нулю і~що це дає для передачі енергії? Як утворюється обертове магнітне поле?
\end{quizlist}
\end{Questions}

\Closesolutionfile{answer}
